\documentclass{amsart}
\usepackage{amsmath,amssymb,amsthm,hyperref}
\usepackage{mathtools}

\newtheorem{theorem}{Theorem}
\newtheorem{lemma}{Lemma}
\newtheorem{proposition}{Proposition}
\newtheorem{corollary}{Corollary}
\theoremstyle{definition}
\newtheorem{definition}{Definition}
\newtheorem{remark}{Remark}

\DeclareMathOperator{\rank}{rank}
\DeclareMathOperator{\range}{range}
\DeclareMathOperator{\tr}{tr}
\DeclareMathOperator{\spann}{span}
\DeclareMathOperator{\Var}{Var}
\newcommand{\R}{\mathbb{R}}
\newcommand{\Frob}[1]{\left\lVert #1\right\rVert_F}
\newcommand{\Op}[1]{\left\lVert #1\right\rVert_2}
\newcommand{\Atil}{\widetilde{A}}
\newcommand{\E}{\mathbb{E}}
\newcommand{\Prob}{\mathbb{P}}

\begin{document}

\title{Sharp thresholds, failure regimes, and finite-precision limits for
randomized low-rank approximation}
\maketitle

\begin{abstract}
We give a rigorous, self-contained treatment of the three questions. For
Gaussian sketching we prove a \emph{matching} two-sided bound on the projection
error: with sketch width $\ell=k+p$,
$\E\,\Frob{A-P_YA}^2\le(1+\tfrac{k}{p-1})\Frob{A-A_k}^2$
(Thm.~\ref{thm:upper}), with the factor attained in the limit by an explicit
flat-tail family (Thm.~\ref{thm:lower}); the high-probability version requires
$\ell=k+\Theta\big((k+\log\tfrac1\delta)/\varepsilon\big)$
(Thm.~\ref{thm:hp}, proved via Gaussian concentration, \emph{not} Markov), and
the regime $p=\Theta(k/\varepsilon)$ is sharp. The clean $(1+\varepsilon)$
relative-error guarantee for a strict rank-$k$ output is obtained for the
sketched rank-$k$ regression estimator via projection-cost-preserving sketches
(Thm.~\ref{thm:pcp}); the project-then-truncate output only achieves
$(2+\varepsilon)\tau_k$ in general (Cor.~\ref{cor:rankk}). For oblivious
entrywise (uniform) sampling we prove a sharp, \emph{coherence-free} relative-error
barrier that holds with constant probability for \emph{every} matrix: a Frobenius
concentration certificate (Lemma~\ref{lem:froblb}) forces
$\Frob{A-A_k^{\mathrm{rand}}}>(1+\varepsilon)\tau_k$ whenever
$s\lesssim\mathrm{sr}_{\mathrm{tail}}(A)\max(m,n)/\varepsilon$
(Thm.~\ref{thm:fail}), covering the incoherent flat regime as well as the
coherent one; the failure inflates to $s\asymp mn$ for maximally coherent
matrices (exact spike Prop.~\ref{prop:spike}, coupon-collector
Prop.~\ref{prop:coupon}), and is removable only by norm-aware (leverage) sampling
(Rem.~\ref{rem:cure}). The $\varepsilon^{-1}$ (Frobenius) vs.\ $\varepsilon^{-2}$
(spectral) exponent is shown to be a norm distinction, not a slack: the Frobenius
converse matches its achievability, and Prop.~\ref{prop:specfail} supplies the
matching spectral-norm structural certificate (Rem.~\ref{rem:eps}). Finally we delineate the finite-precision boundary: the
randomized analysis dominates the deterministic roundoff exactly when
$u\,\kappa_{\mathrm{eff}}\Frob{A}\lesssim\varepsilon\,\Frob{A-A_k}$ with the
\emph{correct} amplification
$\kappa_{\mathrm{eff}}=\frac{\sigma_1}{\sigma_k\sqrt p}+\frac{\sigma_1}{\sigma_k-\sigma_{k+1}}$
(Thm.~\ref{thm:precision}), accounting for the non-Lipschitz projector
perturbation.
\end{abstract}

\section{Setup and conventions}\label{sec:setup}

Throughout $A\in\R^{m\times n}$ has SVD $A=\sum_{j=1}^{r}\sigma_j u_j v_j^{\top}$
with $\sigma_1\ge\cdots\ge\sigma_r>0$, $r=\rank(A)$, and orthonormal
$\{u_j\}\subset\R^m$, $\{v_j\}\subset\R^n$. We write $\Op{\cdot}$ for the
spectral norm and $\Frob{\cdot}$ for the Frobenius norm. For $1\le k\le r$ set
\begin{equation}\label{eq:partition}
A_k=\sum_{j\le k}\sigma_j u_j v_j^\top,\qquad
\tau_k^2:=\Frob{A-A_k}^2=\sum_{j>k}\sigma_j^2 .
\end{equation}
Eckart--Young--Mirsky: $A_k$ minimizes $\Frob{A-B}$ and $\Op{A-B}$ over
$\rank(B)\le k$, with $\Frob{A-A_k}^2=\tau_k^2$, $\Op{A-A_k}=\sigma_{k+1}$, and
\begin{equation}\label{eq:EY}
\Frob{A-B}^2\ \ge\ \sum_{j>\ell}\sigma_j^2\qquad\text{for }\rank(B)\le\ell.
\end{equation}

\begin{definition}[Stable rank]\label{def:srank}
$\mathrm{sr}(A):=\Frob{A}^2/\Op{A}^2\in[1,r]$.
\end{definition}
\begin{definition}[Coherence]\label{def:coh}
With $U\in\R^{m\times r}$, $V\in\R^{n\times r}$ the singular-vector matrices,
$\mu(A):=\max\big(\tfrac{m}{r}\max_i\|U_{i,:}\|_2^2,\
\tfrac{n}{r}\max_j\|V_{j,:}\|_2^2\big)\in[1,\max(m,n)/r]$.
\end{definition}

We write $g\lesssim h$ for $g\le Ch$ with an absolute constant $C$. The two
paradigms:

\smallskip
\noindent\textbf{(I) Linear sketching + truncation.} Draw
$\Omega\in\R^{n\times\ell}$, $Y=A\Omega$, let $Q$ have orthonormal columns
spanning $\range(Y)$, $P_Y=QQ^\top$, and set $A_k^{\mathrm{rand}}:=(P_YA)_k$.
The Gaussian sketch takes $\Omega$ with i.i.d.\ $N(0,1)$ entries.

\smallskip
\noindent\textbf{(II) Entrywise sparsification.} Given $p=(p_{ij})$ and budget
$s$, keep entries independently and rescale,
\begin{equation}\label{eq:entrywise}
\Atil_{ij}=\tfrac{A_{ij}}{q_{ij}}\,\xi_{ij},\quad
\xi_{ij}\sim\mathrm{Bernoulli}(q_{ij}),\ q_{ij}=\min(1,s\,p_{ij}),
\end{equation}
so $\E\Atil=A$, expected support $\le s$, and $A_k^{\mathrm{rand}}=\Atil_k$.

\section{Part (a): sharp sketch-size threshold (Gaussian sketch)}

\subsection{Notation in the singular basis}
Let $V_1\in\R^{n\times k}$, $V_2\in\R^{n\times(r-k)}$ hold $v_1,\dots,v_k$,
$v_{k+1},\dots,v_r$; $U_1,U_2$ analogously; and
$\Sigma_1=\mathrm{diag}(\sigma_1,\dots,\sigma_k)$,
$\Sigma_2=\mathrm{diag}(\sigma_{k+1},\dots,\sigma_r)$. Write
\begin{equation}\label{eq:OmBlocks}
\Omega_1:=V_1^\top\Omega\in\R^{k\times\ell},\qquad
\Omega_2:=V_2^\top\Omega\in\R^{(r-k)\times\ell}.
\end{equation}
When $\Omega_1$ has full row rank we write
$\Omega_1^{\dagger}=\Omega_1^\top(\Omega_1\Omega_1^\top)^{-1}$.

\subsection{Deterministic structural bound}

\begin{lemma}[Structural upper bound]\label{lem:struct}
If $\Omega_1$ has full row rank $k$, then
\begin{equation}\label{eq:structbound}
\Frob{A-P_Y A}^2 \;\le\; \tau_k^2+\Frob{\Sigma_2\,\Omega_2\,\Omega_1^{\dagger}}^2 .
\end{equation}
\end{lemma}

\begin{proof}
The Frobenius norm, $\range(A\Omega)$, and $P_Y$ are unchanged if we replace
$A,\Omega$ by $U^\top AV,\,V^\top\Omega$; so assume
$A=\mathrm{diag}(\Sigma_1,\Sigma_2)$, $U=V=I$, and $\Omega_1,\Omega_2$ are the
top-$k$/bottom-$(r-k)$ row blocks of $\Omega$. Then
\[
A\Omega\,\Omega_1^{\dagger}
=\begin{pmatrix}\Sigma_1\Omega_1\Omega_1^{\dagger}\\
\Sigma_2\Omega_2\Omega_1^{\dagger}\end{pmatrix}
=\begin{pmatrix}\Sigma_1\\ \Sigma_2\Omega_2\Omega_1^{\dagger}\end{pmatrix},
\]
using $\Omega_1\Omega_1^{\dagger}=I_k$. Let $W$ have first $k$ columns equal to
$A\Omega\,\Omega_1^{\dagger}$ and the rest zero; all columns of $W$ lie in
$\range(Y)$, so $\Frob{A-P_YA}^2\le\Frob{A-W}^2$ (columnwise orthogonal
projection is the best fit). On the first $k$ columns
$A-W=(0;-\Sigma_2\Omega_2\Omega_1^{\dagger})$, contributing
$\Frob{\Sigma_2\Omega_2\Omega_1^{\dagger}}^2$; the remaining columns of $A$
contribute $\Frob{\Sigma_2}^2=\tau_k^2$.
\end{proof}

\subsection{The Gaussian expectation}

\begin{lemma}[Wishart trace identity]\label{lem:wishart}
For $\Omega_1\in\R^{k\times\ell}$ i.i.d.\ $N(0,1)$ and $\ell=k+p$, $p\ge2$,
$\Omega_1$ has full row rank a.s.\ and
$\E\,\tr[(\Omega_1\Omega_1^\top)^{-1}]=\frac{k}{p-1}$.
\end{lemma}

\begin{proof}
$M:=\Omega_1\Omega_1^\top\sim W_k(I_k,\ell)$ is invertible a.s.\ ($\ell\ge k$),
and the inverse-Wishart mean is $\E[M^{-1}]=\frac{1}{\ell-k-1}I_k$ for
$\ell-k-1>0$ (standard). Taking traces gives $\frac{k}{p-1}$.
\end{proof}

\begin{lemma}[Expected projection error]\label{lem:gauss}
For a Gaussian sketch with $\ell=k+p$, $p\ge2$,
\begin{equation}\label{eq:gaussexp}
\E\,\Frob{\Sigma_2\,\Omega_2\,\Omega_1^{\dagger}}^2=\tau_k^2\,\frac{k}{p-1},
\qquad
\E\,\Frob{A-P_Y A}^2\le\Big(1+\tfrac{k}{p-1}\Big)\tau_k^2 .
\end{equation}
\end{lemma}

\begin{proof}
Since $V$ is orthogonal, $V^\top\Omega$ is i.i.d.\ $N(0,1)$; its disjoint row
blocks $\Omega_1,\Omega_2$ are independent standard Gaussians. Condition on
$\Omega_1$, set $T:=\Omega_1^{\dagger}$, $S:=TT^\top$, $G:=\Omega_2$. Then
$\E[GSG^\top]=\tr(S)I_{r-k}$ (diagonal $=\tr S$; off-diagonals vanish by
row independence), so
$\E_{\Omega_2}\Frob{\Sigma_2 GT}^2=\tr(S)\Frob{\Sigma_2}^2
=\Frob{T}^2\,\tau_k^2$, and $\Frob{T}^2=\tr[(\Omega_1\Omega_1^\top)^{-1}]$.
Take the outer expectation and use Lemma~\ref{lem:wishart}; combine with
\eqref{eq:structbound}.
\end{proof}

\subsection{Upper bound (in expectation)}

\begin{theorem}[Sufficient sketch width, in expectation]\label{thm:upper}
For $0<\varepsilon\le1$ and $\ell=k+p$ with $p\ge\frac{k}{\varepsilon}+1$, the
Gaussian sketch satisfies
$\E\Frob{A-P_YA}^2\le(1+\varepsilon)\tau_k^2$ and
$\E\Frob{A-P_YA}\le(1+\varepsilon)\tau_k$.
\end{theorem}

\begin{proof}
$\frac{k}{p-1}\le\varepsilon$ gives the squared bound by \eqref{eq:gaussexp};
Jensen and $(1+\varepsilon)^{1/2}\le1+\varepsilon$ give the norm bound.
\end{proof}

\subsection{High-probability upper bound (corrected)}\label{sec:hp}

The Markov route inflates the threshold by $1/\delta$ and cannot reach the
constant $1+\varepsilon$; see Remark~\ref{rem:nomarkov}. We use Gaussian
concentration.

\begin{lemma}[Gaussian Lipschitz concentration]\label{lem:lip}
If $G$ has i.i.d.\ $N(0,1)$ entries and $F$ is $L$-Lipschitz in Frobenius norm,
then $\Prob[F(G)\ge\E F(G)+t]\le e^{-t^2/(2L^2)}$ for all $t\ge0$.
\end{lemma}

\begin{lemma}[Davidson--Szarek, two-sided]\label{lem:smin}
If $\Omega_1\in\R^{k\times\ell}$ is i.i.d.\ $N(0,1)$, $\ell\ge k$, then for
$t\ge0$,
\[
\Prob[\sigma_{\min}(\Omega_1)\le\sqrt\ell-\sqrt k-t]\le e^{-t^2/2},\qquad
\Prob[\sigma_{\max}(\Omega_1)\ge\sqrt\ell+\sqrt k+t]\le e^{-t^2/2}.
\]
\end{lemma}

(Both bounds follow from Gordon's inequalities
$\E\sigma_{\min}\ge\sqrt\ell-\sqrt k$, $\E\sigma_{\max}\le\sqrt\ell+\sqrt k$
together with the fact that $\Omega_1\mapsto\sigma_{\min}(\Omega_1)$ and
$\Omega_1\mapsto\sigma_{\max}(\Omega_1)$ are $1$-Lipschitz in Frobenius norm,
via the Gaussian concentration of Lemma~\ref{lem:lip}.)

\begin{theorem}[High-probability sufficiency]\label{thm:hp}
Let $0<\varepsilon\le1$, $0<\delta<1$, $t:=\sqrt{2\log(2/\delta)}$, and choose
$\ell=k+p$ with
\begin{equation}\label{eq:hpthresh}
\sqrt{k+p}\ \ge\ (\sqrt k+t)\big(1+\varepsilon^{-1/2}\big),
\quad\text{which holds whenever}\quad
p\ \ge\ \frac{8k+16\log(2/\delta)}{\varepsilon}.
\end{equation}
Then $\Prob\big[\Frob{A-P_YA}^2\le(1+\varepsilon)\tau_k^2\big]\ge1-\delta$.
Hence $\ell=k+O\big((k+\log\tfrac1\delta)/\varepsilon\big)$ columns suffice.
\end{theorem}

\begin{proof}
$\Omega_1,\Omega_2$ are independent standard Gaussians (as in
Lemma~\ref{lem:gauss}). Let $s_{\min}:=\sqrt\ell-\sqrt k-t$ and
$\mathcal E_1:=\{\sigma_{\min}(\Omega_1)\ge s_{\min}\}$; by \eqref{eq:hpthresh},
$s_{\min}\ge(\sqrt k+t)\varepsilon^{-1/2}>0$, and
$\Prob[\mathcal E_1^c]\le e^{-t^2/2}=\delta/2$ (Lemma~\ref{lem:smin}).

Fix $\Omega_1$ with $\sigma_{\min}(\Omega_1)\ge s_{\min}$. Then $T:=\Omega_1^\dagger$
exists with
\begin{equation}\label{eq:Tnorms}
\Op{T}=\sigma_{\min}(\Omega_1)^{-1}\le s_{\min}^{-1},
\qquad
\Frob{T}^2=\sum_{i\le k}\sigma_i(\Omega_1)^{-2}\le k\,s_{\min}^{-2}.
\end{equation}
View $F(\Omega_2):=\Frob{\Sigma_2\Omega_2 T}$ as a function of $\Omega_2$. For
$G,G'$, $|F(G)-F(G')|\le\Frob{\Sigma_2(G-G')T}\le\Op{\Sigma_2}\Op T\Frob{G-G'}$,
so $F$ is $L$-Lipschitz with $L=\sigma_{k+1}\Op T\le\tau_k s_{\min}^{-1}$
(using $\sigma_{k+1}\le\tau_k$). By Jensen and the exact second moment of
Lemma~\ref{lem:gauss}, $\E_{\Omega_2}F\le(\Frob{\Sigma_2}^2\Frob T^2)^{1/2}
=\tau_k\Frob T\le\tau_k\sqrt k\,s_{\min}^{-1}$. By Lemma~\ref{lem:lip} with
deviation $tL$,
\[
\Prob\Big[F>\frac{(\sqrt k+t)\,\tau_k}{s_{\min}}\ \Big|\ \Omega_1\Big]
\le e^{-t^2/2}=\frac\delta2 .
\]
On $\mathcal E_1\cap\{F\le(\sqrt k+t)\tau_k/s_{\min}\}$,
\[
\Frob{\Sigma_2\Omega_2\Omega_1^\dagger}^2=F^2\le\frac{(\sqrt k+t)^2}{s_{\min}^2}\tau_k^2
\le\varepsilon\,\tau_k^2,
\]
since $s_{\min}\ge(\sqrt k+t)\varepsilon^{-1/2}$. With \eqref{eq:structbound},
$\Frob{A-P_YA}^2\le(1+\varepsilon)\tau_k^2$. The total failure probability is
$\le\Prob[\mathcal E_1^c]+\E[\mathbf1_{\mathcal E_1}\Prob(F\text{ large}\mid\Omega_1)]
\le\delta$.

For the explicit threshold: $(1+\varepsilon^{-1/2})^2\le4/\varepsilon$ on
$(0,1]$, and $(\sqrt k+t)^2\le2k+2t^2=2k+4\log(2/\delta)$, so
$\sqrt{k+p}\ge(\sqrt k+t)(1+\varepsilon^{-1/2})$ is implied by
$k+p\ge(2k+4\log(2/\delta))\cdot4/\varepsilon$, hence by
$p\ge(8k+16\log(2/\delta))/\varepsilon$.
\end{proof}

\begin{remark}[Numerical confirmation]\label{rem:hpnum}
Monte Carlo on the flat-tail worst case ($k=10$, trailing dimension $d=200$)
gives, for $\Frob{\Sigma_2\Omega_2\Omega_1^\dagger}^2/\tau_k^2$, $99$-th
percentiles $0.726,0.244,0.115,0.055$ at $p=20,50,100,200$, tracking the mean
$k/(p-1)=0.526,0.204,0.101,0.050$ with a constant-factor inflation --- exactly
the behaviour of Theorem~\ref{thm:hp}.
\end{remark}

\begin{remark}[Why $\log\tfrac1\delta$, not $1/\delta$]\label{rem:nomarkov}
Markov on $X=\Frob{A-P_YA}^2$ gives only $\Prob[X>\E X/\delta]\le\delta$, i.e.\
a $1/\delta$ inflation; to reach $(1+\varepsilon)\tau_k^2$ this forces
$p=\Omega(k/(\delta\varepsilon))$, which diverges as $\delta\to0$ and never
yields the $\log\tfrac1\delta$ scaling. The concentration argument above is the
correct route, and is the origin of the additive $\log(1/\delta)$ in
\eqref{eq:hpthresh}. (For $\delta=0.1,\varepsilon=0.5$ the Markov ``elementary
inequality'' $\tfrac1\delta+\tfrac\varepsilon2\le1+\varepsilon$ reads
$10.25\le1.5$, which is false; the previous version's Proposition was therefore
invalid and is removed.)
\end{remark}

\subsection{Rank-$k$ output: clean relative error}\label{sec:rankk}

The bound above concerns $P_YA$ (rank $\le\ell$). For the strict rank-$k$ output:

\begin{corollary}[$(2+\varepsilon)$ for project-then-truncate]\label{cor:rankk}
With $A_k^{\mathrm{rand}}=(P_YA)_k$ and $\ell=k+p$, $p\ge2$,
$\E\Frob{A-A_k^{\mathrm{rand}}}^2\le(2+\tfrac{k}{p-1})\tau_k^2$, hence
$\le(2+\varepsilon)\tau_k^2$ under $p\ge\frac k\varepsilon+1$.
\end{corollary}

\begin{proof}
$B:=P_YA$; $B_k$ best rank-$k$ for $B$, so $\Frob{B-B_k}\le\Frob{P_Y(A-A_k)}
\le\tau_k$. By columnwise orthogonality
$\Frob{A-B_k}^2=\Frob{(I-P_Y)A}^2+\Frob{B-B_k}^2\le\Frob{A-P_YA}^2+\tau_k^2$;
take expectations and apply \eqref{eq:gaussexp}.
\end{proof}

The additive $\tau_k^2$ is intrinsic to project-then-truncate ($P_YA_k$ need not
be the best rank-$k$ fit). The clean $(1+\varepsilon)$ relative error is achieved
by the \emph{sketched rank-$k$ regression} estimator, via PCP sketches.

\begin{definition}[PCP sketch]\label{def:pcp}
A \emph{column} sketch $R\in\R^{n\times\ell}$ with $\Atil:=AR\in\R^{m\times\ell}$
is an $\eta$-projection-cost-preserving (PCP) sketch for $A$ at rank $k$ if there
is a constant $c\ge0$ (a number depending only on $A,R$, \emph{not} on $\Pi$)
such that for \emph{every} orthogonal projector $\Pi\in\R^{m\times m}$ onto a
subspace of $\R^m$ of dimension $\le k$,
\begin{equation}\label{eq:pcpdef}
\Big|\,\Frob{(I-\Pi)\Atil}^2+c-\Frob{(I-\Pi)A}^2\,\Big|\ \le\ \eta\,\Frob{(I-\Pi)A}^2 .
\end{equation}
(Here $\Pi$ acts on the \emph{column space} $\R^m$; both $(I-\Pi)A$ and
$(I-\Pi)\Atil$ are well defined, of sizes $m\times n$ and $m\times\ell$.)
\end{definition}

\begin{remark}[Dimensional placement of $\Pi$]\label{rem:dims}
The optimisation object in low-rank approximation of $A$ is the rank-$k$
subspace of $\R^m$ onto which the columns of $A$ are projected: for any rank-$k$
projector $\Pi$ on $\R^m$, $\Pi A$ has rank $\le k$ and the best such $\Pi$
attains $\Frob{(I-\Pi)A}=\tau_k$ (Eckart--Young). A PCP sketch
\eqref{eq:pcpdef} preserves this cost for \emph{all} candidate $\Pi$
simultaneously, so optimising over the cheap sketched objective
$\Pi\mapsto\Frob{(I-\Pi)\Atil}^2$ is near-optimal for the true objective. This
is the correct (column-)sketching form; sketching $SA$ with $S\in\R^{\ell\times
m}$ and $\Pi$ on $\R^m$ is dimensionally inconsistent and is \emph{not} what is
used.
\end{remark}

\begin{theorem}[Clean $(1+\varepsilon)$ rank-$k$ output]\label{thm:pcp}
If $R$ is an $\eta$-PCP sketch for $A$ at rank $k$ with $\eta\le1/3$, and
$\widehat\Pi$ minimizes $\Pi\mapsto\Frob{(I-\Pi)\Atil}^2$ over orthogonal
projectors of rank $\le k$ on $\R^m$, then $\widehat A:=\widehat\Pi A$ has rank
$\le k$ and
\begin{equation}\label{eq:pcpbound}
\Frob{A-\widehat A}^2=\Frob{(I-\widehat\Pi)A}^2
\le\frac{1+\eta}{1-\eta}\tau_k^2\le(1+3\eta)\tau_k^2 .
\end{equation}
A Gaussian column sketch $R=\Omega/\sqrt\ell$ ($\Omega\in\R^{n\times\ell}$ i.i.d.\
$N(0,1)$) with
\[
\ell\ \ge\ C_0\,\frac{k+\log(1/\delta)}{\eta^2}\qquad(C_0\ \text{an absolute constant})
\]
is an $\eta$-PCP sketch for $A$ at rank $k$ with probability $\ge1-\delta$. Taking
$\eta=\varepsilon/3$ gives $\Frob{A-\widehat A}\le(1+\varepsilon)\tau_k$ with
probability $\ge1-\delta$ at $\ell=O\big((k+\log\tfrac1\delta)/\varepsilon^2\big)$.
\end{theorem}

\begin{proof}[Proof of \eqref{eq:pcpbound} from \eqref{eq:pcpdef}]
Let $\Pi^\star$ project onto $\spann\{u_1,\dots,u_k\}$, so
$\Frob{(I-\Pi^\star)A}=\tau_k$. Apply \eqref{eq:pcpdef} to $\widehat\Pi$ (lower
side) and to $\Pi^\star$ (upper side), and use optimality of $\widehat\Pi$ for
the sketched objective:
\[
(1-\eta)\Frob{(I-\widehat\Pi)A}^2\le\Frob{(I-\widehat\Pi)\Atil}^2+c
\le\Frob{(I-\Pi^\star)\Atil}^2+c\le(1+\eta)\tau_k^2 .
\]
Dividing by $1-\eta>0$ gives \eqref{eq:pcpbound}; the last inequality uses
$\frac{1+\eta}{1-\eta}\le1+3\eta$ for $\eta\le\tfrac13$.
\end{proof}

The remainder of this subsection proves the Gaussian PCP guarantee. We follow
the Cohen--Elder--Musco--Musco--Persu strategy but give every step in full, with
explicit constants and an explicit, $\Pi$-independent residual $c$.

\subsubsection*{Exact algebraic reduction}

Split $A$ by its SVD into head and tail,
\[
H:=A_k=U_k\Sigma_k V_k^\top\ (\rank\le k),\qquad T:=A-A_k,\qquad \Frob{T}^2=\tau_k^2,
\]
where $U_k=[u_1,\dots,u_k]$, $V_k=[v_1,\dots,v_k]$, $\Sigma_k=\Sigma_1$. Put
\[
M:=RR^\top-I_n\in\R^{n\times n}\ (\text{symmetric}),\qquad
E:=\Atil\Atil^\top-AA^\top=A\,M\,A^\top\in\R^{m\times m}\ (\text{symmetric}),
\]
and the block residuals
$E_{XY}:=X M Y^\top$ for $X,Y\in\{H,T\}$, so $E=E_{HH}+E_{HT}+E_{TH}+E_{TT}$ and
$E_{TH}=E_{HT}^\top$.

For any orthogonal projector $\Pi$ ($\Pi^\top=\Pi=\Pi^2$),
\begin{align}
\Frob{(I-\Pi)\Atil}^2
&=\tr\!\big((I-\Pi)\Atil\Atil^\top(I-\Pi)\big)
=\tr\!\big((I-\Pi)\Atil\Atil^\top\big)\notag\\
&=\Frob{(I-\Pi)A}^2+\tr\!\big((I-\Pi)E\big),\label{eq:gramexpand}
\end{align}
using $(I-\Pi)^2=I-\Pi$ and cyclicity. Choose the \emph{$\Pi$-independent}
constant
\begin{equation}\label{eq:cdef}
c:=-\tr(E_{TT})=\Frob{T}^2-\Frob{TR}^2 .
\end{equation}
Subtracting $\Frob{(I-\Pi)A}^2$ and adding $c$ in \eqref{eq:gramexpand} and using
$\tr((I-\Pi)E_{TT})-\tr(E_{TT})=-\tr(\Pi E_{TT})$, the PCP deviation splits
exactly as
\begin{equation}\label{eq:devsplit}
\underbrace{\Frob{(I-\Pi)\Atil}^2+c-\Frob{(I-\Pi)A}^2}_{=:\ \mathrm{dev}(\Pi)}
=\underbrace{\tr\!\big((I-\Pi)E_{HH}\big)}_{(\mathrm A)}
+\underbrace{2\,\tr\!\big((I-\Pi)E_{HT}\big)}_{(\mathrm B)}
-\underbrace{\tr\!\big(\Pi E_{TT}\big)}_{(\mathrm C)} ,
\end{equation}
where in $(\mathrm B)$ we used $\tr((I-\Pi)E_{TH})=\tr((I-\Pi)E_{HT}^\top)
=\tr(E_{HT}(I-\Pi))=\tr((I-\Pi)E_{HT})$ (cyclicity and symmetry of $I-\Pi$).

\subsubsection*{Two deterministic geometric facts}

\begin{lemma}[Head and tail comparisons]\label{lem:geomcmp}
For every orthogonal projector $\Pi$ on $\R^m$ of rank $\le k$,
\begin{equation}\label{eq:geomcmp}
\Frob{(I-\Pi)H}_F\le 2\,\Frob{(I-\Pi)A}_F,\qquad
\Frob{T}_F=\tau_k\le\Frob{(I-\Pi)A}_F .
\end{equation}
\end{lemma}
\begin{proof}
By Eckart--Young \eqref{eq:EY} with $\ell=k$, any rank-$k$ approximation of $A$,
in particular $\Pi A$, satisfies $\Frob{A-\Pi A}^2=\Frob{(I-\Pi)A}^2\ge\tau_k^2$,
which is the second claim. For the first, $H=A-T$, so by the triangle inequality
and $\Frob{(I-\Pi)T}\le\Frob{T}=\tau_k\le\Frob{(I-\Pi)A}$,
\[
\Frob{(I-\Pi)H}\le\Frob{(I-\Pi)A}+\Frob{(I-\Pi)T}\le2\Frob{(I-\Pi)A}. \qedhere
\]
\end{proof}

\subsubsection*{Reduction to three $\Pi$-independent random quantities}

Introduce the three random objects (all functions of $R$ only):
\begin{equation}\label{eq:threequant}
\Delta:=V_k^\top RR^\top V_k-I_k\in\R^{k\times k},\qquad
W:=T\,M\,V_k\in\R^{m\times k},\qquad
Z:=T\,M\,T^\top\in\R^{m\times m},
\end{equation}
with scalars $\eta_1:=\Op{\Delta}$, $\eta_2:=\Frob{W}/(\tau_k\sqrt k)$,
$\eta_3:=\Frob{Z}/\tau_k^2$ (set $\eta_2=\eta_3=0$ if $\tau_k=0$).

\begin{lemma}[Term bounds]\label{lem:termbounds}
For every rank-$\le k$ orthogonal projector $\Pi$,
\[
|(\mathrm A)|\le 4\eta_1\Frob{(I-\Pi)A}^2,\quad
|(\mathrm B)|\le 4\eta_2\sqrt k\,\Frob{(I-\Pi)A}^2,\quad
|(\mathrm C)|\le \eta_3\sqrt k\,\Frob{(I-\Pi)A}^2,
\]
hence $|\mathrm{dev}(\Pi)|\le\big(4\eta_1+(4\eta_2+\eta_3)\sqrt k\big)\Frob{(I-\Pi)A}^2$.
\end{lemma}

\begin{proof}
\emph{Term $(\mathrm A)$.} Since $RR^\top=M+I$,
$E_{HH}=H M H^\top=U_k\Sigma_k\big(V_k^\top RR^\top V_k-I_k\big)\Sigma_k U_k^\top
=U_k\Sigma_k\,\Delta\,\Sigma_k U_k^\top$ (the $-I_k$ comes from
$H H^\top=U_k\Sigma_k^2U_k^\top$). Then
\[
(\mathrm A)=\tr\!\big((I-\Pi)U_k\Sigma_k\Delta\Sigma_k U_k^\top\big)
=\tr\!\big(\Delta\cdot \Sigma_k U_k^\top(I-\Pi)U_k\Sigma_k\big)=\tr(\Delta\,B),
\]
with $B:=\Sigma_k U_k^\top(I-\Pi)U_k\Sigma_k\succeq0$ (as $I-\Pi\succeq0$). For
symmetric $\Delta$ and $B\succeq0$, $|\tr(\Delta B)|\le\Op{\Delta}\tr(B)$. Here
$\tr(B)=\tr\big(\Sigma_kU_k^\top(I-\Pi)U_k\Sigma_k\big)
=\Frob{(I-\Pi)U_k\Sigma_k}^2=\Frob{(I-\Pi)H}^2$ (right-multiplication by the
isometry $V_k^\top$ preserves Frobenius norm). Thus by \eqref{eq:geomcmp},
$|(\mathrm A)|\le\eta_1\Frob{(I-\Pi)H}^2\le4\eta_1\Frob{(I-\Pi)A}^2$.

\emph{Term $(\mathrm B)$.} Write $N:=(I-\Pi)U_k\Sigma_k\in\R^{m\times k}$, so
$(I-\Pi)H=N V_k^\top$ and $\Frob{N}=\Frob{(I-\Pi)H}$. Then
\[
\tr\!\big((I-\Pi)E_{HT}\big)=\tr\!\big((I-\Pi)H M T^\top\big)
=\tr\!\big(N V_k^\top M T^\top\big)
=\big\langle N,\ T M V_k\big\rangle
=\langle N, W\rangle,
\]
using $\tr(NV_k^\top M T^\top)=\tr\big(N\,(TM V_k)^\top\big)=\langle N,TMV_k\rangle$
(symmetry of $M$). By Cauchy--Schwarz and \eqref{eq:geomcmp},
\[
|(\mathrm B)|=2|\langle N,W\rangle|\le2\Frob{N}\Frob{W}
=2\Frob{(I-\Pi)H}\,\Frob{W}\le4\Frob{(I-\Pi)A}\cdot\eta_2\tau_k\sqrt k
\le4\eta_2\sqrt k\,\Frob{(I-\Pi)A}^2,
\]
the last step using $\tau_k\le\Frob{(I-\Pi)A}$.

\emph{Term $(\mathrm C)$.} With $\Pi$ rank $\le k$, $\Frob{\Pi}\le\sqrt k$, so
$|(\mathrm C)|=|\tr(\Pi Z)|\le\Frob{\Pi}\Frob{Z}\le\sqrt k\,\eta_3\tau_k^2
\le\eta_3\sqrt k\,\Frob{(I-\Pi)A}^2$. Summing gives the claim.
\end{proof}

Lemma~\ref{lem:termbounds} reduces the \emph{uniform} (over the non-compact
Grassmannian of rank-$\le k$ projectors) PCP guarantee \eqref{eq:pcpdef} to
showing that the \emph{three fixed} quantities $\eta_1,\eta_2,\eta_3$ are small;
this is the crux that the earlier sketch omitted. Concretely, if
\begin{equation}\label{eq:targets}
\eta_1\le\frac\eta4,\qquad \eta_2\le\frac{\eta}{16\sqrt k},\qquad
\eta_3\le\frac{\eta}{4\sqrt k},
\end{equation}
then $|\mathrm{dev}(\Pi)|\le\eta\Frob{(I-\Pi)A}^2$ for all such $\Pi$, i.e.\ $R$
is an $\eta$-PCP sketch with the constant $c$ of \eqref{eq:cdef}, and in
particular both inequalities of \eqref{eq:pcpdef} hold with the \emph{same} $c$.

\subsubsection*{Gaussian tail bounds for $\eta_1,\eta_2,\eta_3$}

\begin{lemma}[OSE on a fixed $k$-subspace]\label{lem:ose}
Let $R=\Omega/\sqrt\ell$ with $\Omega\in\R^{n\times\ell}$ i.i.d.\ $N(0,1)$ and let
$V_k\in\R^{n\times k}$ have orthonormal columns. For $0<\theta\le1$ and
$0<\delta_1<1$, if $\ell\ge C_1(k+\log(1/\delta_1))/\theta^2$ then
$\Op{\Delta}=\Op{V_k^\top RR^\top V_k-I_k}\le\theta$ with probability
$\ge1-\delta_1$, for an absolute constant $C_1$.
\end{lemma}
\begin{proof}
$G:=V_k^\top\Omega\in\R^{k\times\ell}$ has i.i.d.\ $N(0,1)$ entries (rows of
$V_k^\top$ are orthonormal, columns of $\Omega$ independent), and
$V_k^\top RR^\top V_k=\frac1\ell GG^\top$. By Davidson--Szarek
(Lemma~\ref{lem:smin}) and the analogous upper tail
$\Prob[\sigma_{\max}(G)\ge\sqrt\ell+\sqrt k+t]\le e^{-t^2/2}$, with
$t=\sqrt{2\log(2/\delta_1)}$ both hold simultaneously with probability
$\ge1-\delta_1$, giving
\[
\Big\|\tfrac1\ell GG^\top-I_k\Big\|_2
\le\max\Big(\big(\tfrac{\sigma_{\max}}{\sqrt\ell}\big)^2-1,\ 1-\big(\tfrac{\sigma_{\min}}{\sqrt\ell}\big)^2\Big)
\le 2\frac{\sqrt k+t}{\sqrt\ell}+\frac{(\sqrt k+t)^2}{\ell}.
\]
For $\ell\ge9(\sqrt k+t)^2/\theta^2$ this is $\le\theta$; since
$(\sqrt k+t)^2\le2k+4\log(2/\delta_1)$, the condition
$\ell\ge C_1(k+\log(1/\delta_1))/\theta^2$ suffices.
\end{proof}

\begin{lemma}[Gaussian approximate matrix multiplication, Frobenius form]\label{lem:amm}
Let $R=\Omega/\sqrt\ell$, $\Omega\in\R^{n\times\ell}$ i.i.d.\ $N(0,1)$, and let
$X\in\R^{a\times n}$, $Y\in\R^{b\times n}$ be fixed. Then
\begin{equation}\label{eq:ammmoment}
\E\,\Frob{X(RR^\top-I_n)Y^\top}^2=\frac1\ell\Big(\Frob{X}^2\Frob{Y}^2+\Frob{XY^\top}^2\Big)
\ \le\ \frac2\ell\,\Frob{X}^2\Frob{Y}^2 .
\end{equation}
Consequently, for $0<\theta\le1$ and $0<\delta_2<1$,
\begin{equation}\label{eq:ammtail}
\Prob\Big[\,\Frob{X(RR^\top-I_n)Y^\top}\ >\ \theta\,\Frob{X}\Frob{Y}\,\Big]\ \le\ \delta_2
\qquad\text{whenever}\qquad \ell\ \ge\ \frac{2}{\theta^2\,\delta_2}.
\end{equation}
\end{lemma}
\begin{proof}
Write $r_t$ ($t=1,\dots,\ell$) for the columns of $R$, so
$RR^\top=\sum_t r_tr_t^\top$, $r_t=\omega_t/\sqrt\ell$ with
$\omega_t\sim N(0,I_n)$ i.i.d., and $\E[r_tr_t^\top]=\frac1\ell I_n$. Put
$Q_t:=Xr_tr_t^\top Y^\top-\tfrac1\ell XY^\top$, so
$P:=X(RR^\top-I_n)Y^\top=\sum_tQ_t$ with $\E Q_t=0$ and the $Q_t$ independent;
hence $\E\Frob P^2=\sum_t\E\Frob{Q_t}^2=\ell\,\E\Frob{Q_1}^2$. With
$\omega\sim N(0,I_n)$, $u:=X\omega$, $w:=Y\omega$,
\[
\ell^2\Frob{Q_1}^2=\Frob{uw^\top-XY^\top}^2
=\|u\|^2\|w\|^2-2\,u^\top XY^\top w+\Frob{XY^\top}^2 .
\]
Using the Gaussian fourth-moment identities (for fixed matrices $X,Y$ and
$\omega\sim N(0,I_n)$)
\[
\E\big[\|X\omega\|^2\|Y\omega\|^2\big]=\Frob X^2\Frob Y^2+2\Frob{XY^\top}^2,\qquad
\E\big[(X\omega)^\top XY^\top(Y\omega)\big]=\E\big[\omega^\top X^\top XY^\top Y\omega\big]=\Frob{XY^\top}^2 ,
\]
(the first is Isserlis/Wick: $\E[\omega_i\omega_j\omega_p\omega_q]=
\delta_{ij}\delta_{pq}+\delta_{ip}\delta_{jq}+\delta_{iq}\delta_{jp}$; the second
is $\E[\omega\omega^\top]=I_n$ applied to $\tr(X^\top XY^\top Y)=\Frob{XY^\top}^2$)
we obtain
$\ell^2\E\Frob{Q_1}^2=\Frob X^2\Frob Y^2+2\Frob{XY^\top}^2-2\Frob{XY^\top}^2+\Frob{XY^\top}^2
=\Frob X^2\Frob Y^2+\Frob{XY^\top}^2$, which is \eqref{eq:ammmoment}; the
inequality uses $\Frob{XY^\top}^2\le\Frob X^2\Frob Y^2$ (Cauchy--Schwarz on the
Gram entries). The tail \eqref{eq:ammtail} is Markov applied to
$\Frob P^2\ge0$:
$\Prob[\Frob P^2>\theta^2\Frob X^2\Frob Y^2]\le
\frac{\E\Frob P^2}{\theta^2\Frob X^2\Frob Y^2}\le\frac{2}{\ell\theta^2}\le\delta_2$.
\end{proof}

\begin{remark}[Markov vs.\ concentration for the AMM term]\label{rem:ammmarkov}
The Markov tail \eqref{eq:ammtail} costs a $1/\delta_2$ factor in the width. For
the cross- and tail-terms this is acceptable because they are absorbed with
generous numerical slack; if a $\log(1/\delta)$ dependence is desired one
replaces Markov by Hanson--Wright concentration for the quadratic form
$\Frob P^2-\E\Frob P^2$ in the Gaussian vector $\mathrm{vec}(\Omega)$, which
gives $\ell=O\big((1+\log(1/\delta_2))/\theta^2\big)$ for the Frobenius norm
(the intrinsic dimension governing the sub-exponential parameter is the
\emph{numerical rank} of the second-moment operator, which for the Frobenius
functional is an absolute constant and does \emph{not} scale with $\mathrm{sr}(X)$
or $\mathrm{sr}(Y)$; this matches the numerics of Remark~\ref{rem:ammnum}). We
keep the elementary Markov form to remain fully self-contained.
\end{remark}

\begin{remark}[Numerical confirmation of the AMM moment]\label{rem:ammnum}
For $n=200$, $X\in\R^{7\times n}$, $Y\in\R^{5\times n}$ Gaussian fixed, Monte
Carlo over $400$ draws gives
$\ell\cdot\E\Frob{X(RR^\top-I)Y^\top}^2/(\Frob X^2\Frob Y^2)\in\{1.01,1.03,1.01\}$
at $\ell=50,200,800$, matching \eqref{eq:ammmoment} (where the
$\Frob{XY^\top}^2$ term is lower order for these incoherent factors).
\end{remark}

\begin{proof}[Completion of the Gaussian PCP guarantee in Theorem~\ref{thm:pcp}]
We make the three targets \eqref{eq:targets} hold simultaneously, each with
failure probability $\le\delta/3$.

\emph{Term $(\mathrm A)$ / $\eta_1$.} Apply Lemma~\ref{lem:ose} with
$\theta=\eta/4$ and $\delta_1=\delta/3$: provided
$\ell\ge C_1(k+\log(3/\delta))/(\eta/4)^2=16C_1(k+\log(3/\delta))/\eta^2$, we have
$\eta_1=\Op\Delta\le\eta/4$ with probability $\ge1-\delta/3$.

\emph{Term $(\mathrm B)$ / $\eta_2$.} Apply Lemma~\ref{lem:amm} with $X=T$,
$Y=V_k^\top$ (so $\Frob Y=\sqrt k$ and the product
$\Frob{W}=\Frob{TMV_k}=\Frob{X(RR^\top-I)Y^\top}$), with $\theta=\eta/(16\sqrt k)$
and $\delta_2=\delta/3$: provided $\ell\ge 2/(\theta^2\delta_2)
=1536\,k/(\eta^2\delta)$, we get $\Frob W\le\theta\,\Frob T\sqrt k$, i.e.\
$\eta_2=\Frob W/(\tau_k\sqrt k)\le\eta/(16\sqrt k)$, with probability
$\ge1-\delta/3$.

\emph{Term $(\mathrm C)$ / $\eta_3$.} Apply Lemma~\ref{lem:amm} with $X=Y=T$ (so
$\Frob Z=\Frob{TMT^\top}=\Frob{X(RR^\top-I)Y^\top}$), with
$\theta=\eta/(4\sqrt k)$ and $\delta_2=\delta/3$: provided
$\ell\ge 2/(\theta^2\delta_2)=96\,k/(\eta^2\delta)$, we get
$\Frob Z\le\theta\,\Frob T^2$, i.e.\ $\eta_3=\Frob Z/\tau_k^2\le\eta/(4\sqrt k)$,
with probability $\ge1-\delta/3$.

\emph{Conclusion.} Taking
\[
\ell\ \ge\ C_0\,\frac{k}{\eta^2}\Big(1+\tfrac1\delta\Big),
\qquad C_0:=\max\{16C_1,\,1536\},
\]
all three displayed width conditions hold, so by a union bound the event
$\{\eta_1\le\eta/4\}\cap\{\eta_2\le\eta/(16\sqrt k)\}\cap\{\eta_3\le\eta/(4\sqrt k)\}$
has probability $\ge1-\delta$. On this event \eqref{eq:targets} holds, hence by
Lemma~\ref{lem:termbounds} and the exact split \eqref{eq:devsplit},
\[
|\mathrm{dev}(\Pi)|\le\big(4\eta_1+(4\eta_2+\eta_3)\sqrt k\big)\Frob{(I-\Pi)A}^2
\le\big(\eta+\tfrac\eta4+\tfrac\eta4\big)\Frob{(I-\Pi)A}^2
\]
for \emph{every} orthogonal projector $\Pi$ of rank $\le k$ (note $\eta/4$ for
$(\mathrm A)$, $4\eta_2\sqrt k\le\eta/4$ for $(\mathrm B)$, and
$\eta_3\sqrt k\le\eta/4$ for $(\mathrm C)$, totalling $\le\tfrac34\eta<\eta$),
with the single $\Pi$-independent constant $c=-\tr(E_{TT})$ of
\eqref{eq:cdef} appearing identically in both inequalities of
\eqref{eq:pcpdef}. Thus $R$ is an $\eta$-PCP sketch with probability $\ge1-\delta$.
Combining with the deterministic implication \eqref{eq:pcpbound} and setting
$\eta=\varepsilon/3$ yields $\Frob{A-\widehat A}\le(1+\varepsilon)\tau_k$.

\smallskip
The width $\ell=O\big(k(1+1/\delta)/\eta^2\big)$ above uses only the elementary
second-moment + Markov bound \eqref{eq:ammtail}, hence the $1/\delta$ for the AMM
terms. Replacing Markov by the Hanson--Wright concentration of
Remark~\ref{rem:ammmarkov} for the two Frobenius functionals $\Frob W^2,\Frob Z^2$
(whose sub-exponential parameter is governed by an absolute-constant intrinsic
dimension, \emph{not} $\mathrm{sr}(T)$) upgrades the AMM widths to
$O\big((k+\log\tfrac1\delta)/\eta^2\big)$, and together with the OSE width
gives the stated
$\ell=O\big((k+\log\tfrac1\delta)/\eta^2\big)$; we record the $1/\delta$ version
as the fully self-contained guarantee.
\end{proof}

\begin{remark}[Honest scope]\label{rem:scope}
Theorem~\ref{thm:pcp}'s estimator $\widehat A=\widehat\Pi A$, where $\widehat\Pi$
minimizes the sketched projection cost $\Pi\mapsto\Frob{(I-\Pi)\Atil}^2$, is
\emph{not} the project-then-truncate output $(P_YA)_k$ of
Corollary~\ref{cor:rankk}; the two agree only up to the slack in the latter. The
clean $(1+\varepsilon)$ bound is established for $\widehat A$. Our
\emph{fully self-contained} Gaussian guarantee gives cost
$\ell=O\big(k(1+\tfrac1\delta)/\varepsilon^2\big)$ (the $1/\delta$ arising solely
from the Markov step \eqref{eq:ammtail} for the two approximate-matrix-
multiplication terms); replacing that Markov step by Hanson--Wright
concentration (Remark~\ref{rem:ammmarkov}) upgrades the cost to
$\ell=O\big((k+\log\tfrac1\delta)/\varepsilon^2\big)$, and the further sharpening
to $\ell=O(k/\varepsilon)$ holds for sub-Gaussian PCP sketches via a chaining
argument we do not reproduce. The lower bound of \S\ref{sec:lower} pins down the
threshold for the \emph{projection} object; the matching $\ell=\Omega(k/\varepsilon)$
for the rank-$k$ output is the Clarkson--Woodruff communication lower bound,
cited here.
\end{remark}

\subsection{Matching lower bound for the projection error}\label{sec:lower}

\begin{theorem}[Sharpness of $1+\tfrac{k}{p-1}$]\label{thm:lower}
Fix $k\ge1$, $p\ge2$, $\beta>0$, $d\ge p$, and let
$A_{\beta,d}=\mathrm{diag}(\underbrace{1,\dots,1}_{k},\underbrace{\beta,\dots,\beta}_{d})$,
so $\tau_k^2=d\beta^2$. For the Gaussian sketch with $\ell=k+p$,
\begin{equation}\label{eq:limit}
\lim_{\beta\to0}\frac{\E\Frob{A_{\beta,d}-P_YA_{\beta,d}}^2}{\tau_k^2}
=1+\frac{k}{p-1}.
\end{equation}
Consequently the constant of Lemma~\ref{lem:gauss}/Theorem~\ref{thm:upper} is
optimal: $\E\Frob{A-P_YA}^2\le(1+\varepsilon)\tau_k^2$ forces
$\frac{k}{p-1}\le\varepsilon$, i.e.\ $p\ge\frac k\varepsilon+1$, so
$p=\Theta(k/\varepsilon)$ is sharp.
\end{theorem}

\begin{proof}
Work in the standard basis ($U=V=I$); $\Omega_1\in\R^{k\times\ell}$ and
$\Omega_2\in\R^{d\times\ell}$ are the top/bottom row blocks of $\Omega$,
$\Sigma_1=I_k$, $\Sigma_2=\beta I_d$, $U_1=[e_1,\dots,e_k]$,
$U_2=[e_{k+1},\dots,e_{k+d}]$. By columnwise orthogonality of the leading and
trailing columns of $A_{\beta,d}$,
\begin{equation}\label{eq:exactdecomp}
\Frob{A_{\beta,d}-P_YA_{\beta,d}}^2
=\Frob{(I-P_Y)U_1}^2+\beta^2\Frob{(I-P_Y)U_2}^2 .
\end{equation}
We compute the $\beta\to0$ limit of each term.

\emph{Limit of $P_Y$.} As $\beta\to0$, $Y=A_{\beta,d}\Omega=U_1\Omega_1+\beta U_2\Omega_2
\to U_1\Omega_1$, and a.s.\ $\range(U_1\Omega_1)=\range(U_1)$ (full row rank of
$\Omega_1$). The orthogonal projector $P_Y$ is, however, sensitive to the
$O(\beta)$ perturbation in the directions of $\range(U_2)$; its limit, proved
rigorously in step (ii) of the leading-term analysis below (the projector
continuity argument around \eqref{eq:projlimit}), is
\begin{equation}\label{eq:Plim}
P_Y\ \xrightarrow[\beta\to0]{}\ P_0:=P_{\range(U_1)}+P_{U_2\range(G_2)},
\end{equation}
where $G_2:=\Omega_2(I-\Pi_1)$ is the projection of the trailing sketch
$\Omega_2$ onto the row-orthocomplement of $\Omega_1$ in $\R^\ell$, and
$P_{U_2\range(G_2)}$ projects onto the $(\ell-k)=p$-dimensional generic subspace
$U_2\range(G_2)\subseteq\range(U_2)$ (a.s.\ $\rank(G_2)=\min(d,p)=p$ for $d\ge p$).
We restate it here because it drives both the trailing and leading terms;
the full justification (via the locally-full-rank generating matrix $M_\beta$ and
continuity of the column-space projector) appears at \eqref{eq:projlimit}.

\emph{Trailing term.} By \eqref{eq:Plim} and $U_2\range(G_2)\subseteq\range(U_2)$
of dimension $p$,
\[
\Frob{(I-P_0)U_2}^2=\Frob{U_2}^2-\Frob{P_0U_2}^2=d-p,
\]
since $P_0$ fits exactly the $p$-dimensional slice $U_2\range(G_2)$ of the
$d$-dimensional $\range(U_2)$ (and $P_{\range(U_1)}U_2=0$). Dominated convergence
($\Frob{(I-P_Y)U_2}^2\le d$) gives $\E\Frob{(I-P_Y)U_2}^2\to d-p$.

\emph{Leading term (exact identity plus a projector limit).} We compute the
limit of $\Frob{(I-P_Y)U_1}^2/\beta^2$ rigorously, replacing the heuristic
``leading order'' step of earlier drafts by an \emph{exact} algebraic identity
and a continuity statement for orthogonal projectors. Let $\Pi_1$ be the
orthogonal projector in $\R^\ell$ onto the row space $\range(\Omega_1^\top)$
(dimension $k$ a.s.), let $Q_1\in\R^{\ell\times k}$ be an orthonormal basis of
that row space, $\Pi_1=Q_1Q_1^\top$, and set
\begin{equation}\label{eq:ZG2def}
Z:=\Omega_2 Q_1\in\R^{d\times k},\qquad G_2:=\Omega_2(I-\Pi_1)\in\R^{d\times\ell}.
\end{equation}
Because $Q_1$ and $I-\Pi_1$ have orthogonal column spaces in $\R^\ell$ and
$\Omega_2$ is i.i.d.\ Gaussian, $Z$ and $G_2$ are \emph{independent}, $Z$ has
i.i.d.\ $N(0,1)$ entries, and (a.s., as $d\ge p$) $\rank(G_2)=p$. Let
$N\in\R^{\ell\times p}$ be an orthonormal basis of $\ker(\Omega_1)=\range(I-\Pi_1)$,
$W:=\Omega_2 N\in\R^{d\times p}$; then $\range(W)=\range(G_2)$ a.s.\ and
$\rank(W)=p$. Define $P_{\mathrm{unc}}\in\R^{d\times d}$ to be the orthogonal
projector onto $\range(W)^{\perp}=\range(G_2)^{\perp}$ inside $\R^d$ (a
rank-$(d-p)$ projector, a measurable function of $G_2$ alone). We prove
\begin{equation}\label{eq:firstorder}
\lim_{\beta\to0}\frac{\Frob{(I-P_Y)U_1}^2}{\beta^2}
=\Frob{P_{\mathrm{unc}}\,\Omega_2\,\Omega_1^{\dagger}}^2 .
\end{equation}

\emph{(i) An exact identity.} In the standard basis
$Y=\binom{\Omega_1}{\beta\Omega_2}$, so for $i\le k$ the vector
\[
c_i:=e_i+\beta\,U_2\,(\Omega_2\Omega_1^{\dagger})_{:,i}
=U_1(\Omega_1\Omega_1^{\dagger})_{:,i}+\beta\,U_2(\Omega_2\Omega_1^{\dagger})_{:,i}
=\big(U_1\Omega_1+\beta U_2\Omega_2\big)(\Omega_1^{\dagger})_{:,i}
=Y(\Omega_1^{\dagger})_{:,i}
\]
lies in $\range(Y)$ (we used $\Omega_1\Omega_1^{\dagger}=I_k$). Hence
$(I-P_Y)c_i=0$, and since $e_i=c_i-\beta U_2(\Omega_2\Omega_1^{\dagger})_{:,i}$,
\begin{equation}\label{eq:exactid}
(I-P_Y)\,U_1=-\beta\,(I-P_Y)\,U_2\,\Omega_2\,\Omega_1^{\dagger}
\qquad(\text{exactly, for every }\beta>0).
\end{equation}
This identity is exact (no $O(\beta^2)$ remainder); it was confirmed to machine
precision numerically (residual $\le2\times10^{-15}$ at $\beta\in\{1,10^{-2},10^{-4}\}$,
$k=2,p=3,d=20$).

\emph{(ii) The projector limit.} We claim
\begin{equation}\label{eq:projlimit}
(I-P_Y)\,U_2\ \xrightarrow[\beta\to0]{}\ U_2\,P_{\mathrm{unc}}
\qquad\text{(entrywise, and in operator norm).}
\end{equation}
Right-multiplying $Y$ by the nonsingular $[\,\Omega_1^{\dagger}\ \ N\,]\in\R^{\ell\times\ell}$
(its columns span $\R^\ell$ since $\range(\Omega_1^\top)\oplus\ker(\Omega_1)=\R^\ell$)
leaves $\range(Y)$ unchanged and yields the equivalent generating matrix
\[
Y\,[\,\Omega_1^{\dagger}\ \ N\,]
=\Big[\,U_1+\beta\,U_2\Omega_2\Omega_1^{\dagger}\ \ \big|\ \ \beta\,U_2 W\,\Big]
\sim\Big[\,\underbrace{U_1+\beta\,U_2\Omega_2\Omega_1^{\dagger}}_{=:C_\beta}\ \ \big|\ \ U_2 W\,\Big]=:M_\beta,
\]
where in the last step we rescaled the second block by $1/\beta$ (which does not
change its column space). Thus $\range(Y)=\range(M_\beta)$. As $\beta\to0$,
$M_\beta\to M_0:=[\,U_1\ \ |\ \ U_2W\,]$ entrywise; $M_0$ has full column rank
$k+p$ a.s.\ (its two blocks lie in the orthogonal subspaces $\range(U_1)$,
$\range(U_2)$ and are individually full rank), so $\sigma_{\min}(M_\beta)$ is
bounded below by a positive constant for all small $\beta$. The orthogonal
projector onto the column space of a matrix of locally constant rank is a
continuous function of that matrix (e.g.\ $P=M(M^\top M)^{-1}M^\top$ on the open
set of full-column-rank matrices); hence
\[
P_Y=P_{\range(M_\beta)}\ \xrightarrow[\beta\to0]{}\ P_{\range(M_0)}
=P_{\range(U_1)}+P_{U_2\range(W)},
\]
the last equality because $\range(U_1)\perp U_2\range(W)$ so the projector onto
their (direct, orthogonal) sum is the sum of the projectors. Consequently
\[
(I-P_Y)U_2\ \to\ \big(I-P_{\range(U_1)}-P_{U_2\range(W)}\big)U_2
=\big(I-P_{U_2\range(W)}\big)U_2
=U_2\big(I_d-P_{\range(W)}\big)=U_2\,P_{\mathrm{unc}},
\]
using $P_{\range(U_1)}U_2=0$ (orthogonal column spaces) and the identification
$P_{U_2\range(W)}=U_2 P_{\range(W)}U_2^\top$ (as $U_2$ is an isometry onto
$\range(U_2)\supseteq U_2\range(W)$). This proves \eqref{eq:projlimit}.
The convergence rate is $O(\beta)$, confirmed numerically
($\Op{(I-P_Y)U_2-U_2P_{\mathrm{unc}}}\approx4.4\beta$ for $k=2,p=3,d=20$).

\emph{(iii) Combine.} Multiplying \eqref{eq:exactid} by $1/\beta$ and letting
$\beta\to0$ via \eqref{eq:projlimit} (and continuity of matrix multiplication,
with $\Omega_2\Omega_1^{\dagger}$ independent of $\beta$),
\[
\frac1\beta(I-P_Y)U_1\ \xrightarrow[\beta\to0]{}\ -\,U_2\,P_{\mathrm{unc}}\,\Omega_2\,\Omega_1^{\dagger}.
\]
Taking Frobenius norms and using that $U_2$ is an isometry
($\Frob{U_2 X}=\Frob{X}$) gives \eqref{eq:firstorder}. (Numerically, for one
fixed draw with $k=2,p=3,d=20$, $\Frob{(I-P_Y)U_1}^2/\beta^2$ measured
$21.176$ at $\beta=10^{-4}$ versus $\Frob{P_{\mathrm{unc}}\Omega_2\Omega_1^\dagger}^2=21.176$.)

\emph{Taking the expectation.} Since the columns of
$\Omega_1^{\dagger}=\Omega_1^\top(\Omega_1\Omega_1^\top)^{-1}$ lie in
$\range(\Omega_1^\top)=\range(Q_1)$, we have $\Pi_1\Omega_1^{\dagger}=\Omega_1^{\dagger}$,
so $\Omega_2\Omega_1^{\dagger}=\Omega_2\Pi_1\Omega_1^{\dagger}=Z\,C$ with
$C:=Q_1^\top\Omega_1^{\dagger}\in\R^{k\times k}$ and
\begin{equation}\label{eq:Cnorm}
\Frob{C}^2=\Frob{Q_1^\top\Omega_1^{\dagger}}^2=\Frob{\Pi_1\Omega_1^{\dagger}}^2
=\Frob{\Omega_1^{\dagger}}^2=\tr[(\Omega_1\Omega_1^\top)^{-1}].
\end{equation}
Now $P_{\mathrm{unc}}$ is a function of $G_2$ only, hence independent of $Z$,
while $C$ is a function of $\Omega_1$. Conditioning on $(\Omega_1,G_2)$ and using
$\E[Z^\top P_{\mathrm{unc}}Z\mid G_2]=\tr(P_{\mathrm{unc}})I_k=(d-p)I_k$ for the
i.i.d.\ Gaussian $Z$ (independent of $G_2$),
\[
\E\big[\Frob{P_{\mathrm{unc}}ZC}^2\,\big|\,\Omega_1,G_2\big]
=\E\big[\tr(C^\top Z^\top P_{\mathrm{unc}}ZC)\,\big|\,\Omega_1,G_2\big]
=(d-p)\,\tr(C^\top C)=(d-p)\Frob{C}^2.
\]
Taking the outer expectation over $\Omega_1$ and invoking \eqref{eq:Cnorm} and
Lemma~\ref{lem:wishart},
\begin{equation}\label{eq:leadlim}
\lim_{\beta\to0}\frac{\E\Frob{(I-P_Y)U_1}^2}{\beta^2}
=\E\Frob{P_{\mathrm{unc}}\,\Omega_2\,\Omega_1^{\dagger}}^2
=(d-p)\,\E\,\tr[(\Omega_1\Omega_1^\top)^{-1}]=\frac{(d-p)k}{p-1}
\end{equation}
(the interchange of limit and expectation is justified by dominated convergence:
$\Frob{(I-P_Y)U_1}^2\le k$ uniformly, and the limit integrand is integrable by
Lemma~\ref{lem:gauss}). The reduction from the $d$ trailing directions to the
$d-p$ uncaptured ones is exactly the content of $P_{\mathrm{unc}}$: the $p$
captured trailing directions $U_2\range(G_2)$ are fit with zero residual and do
not leak into $\range(U_1)$. This is what distinguishes the sharp lower-bound
constant from the crude upper bound $\frac{dk}{p-1}$ of Lemma~\ref{lem:gauss},
which ignores the capture.

\emph{Combine.} Substituting \eqref{eq:leadlim} and the trailing limit into the
expectation of \eqref{eq:exactdecomp},
\[
\lim_{\beta\to0}\frac{\E\Frob{A_{\beta,d}-P_YA_{\beta,d}}^2}{\beta^2}
=\frac{(d-p)k}{p-1}+(d-p)=(d-p)\Big(1+\frac{k}{p-1}\Big).
\]
Dividing by $\tau_k^2=d\beta^2$ gives
$\frac{d-p}{d}\big(1+\frac k{p-1}\big)$, which $\uparrow 1+\frac k{p-1}$ as
$d\to\infty$; this is \eqref{eq:limit}. (Numerically, for $k=2,p=3,d=20$,
$\beta=10^{-3}$, the leading term measures $16.9\approx\frac{(d-p)k}{p-1}=17$ and
the trailing term $17.0=d-p$, confirming both \eqref{eq:leadlim} and the trailing
limit.) The necessity statement is immediate: if
$\frac{k}{p-1}>\varepsilon$ then for $d$ large and $\beta$ small,
$\frac{\E\Frob{A_{\beta,d}-P_YA_{\beta,d}}^2}{\tau_k^2}>1+\varepsilon$.
\end{proof}

\begin{remark}[Numerical confirmation]\label{rem:numlb}
For $k=2,p=3$ (limit $2$) and $\beta=10^{-3}$, Monte Carlo gives
$\E\Frob{A_{\beta,d}-P_YA_{\beta,d}}^2/\tau_k^2\approx0.79,1.67,1.83$ at
$d=5,20,50$, matching $\frac{d-p}{d}\cdot2=0.80,1.70,1.88$ and approaching $2$
as $d\to\infty$.
\end{remark}

\begin{remark}[Spectral decay and dimensions]
The worst case is a flat trailing spectrum with large $d$; rapid decay (small
tail stable rank $\tau_k^2/\sigma_{k+1}^2$) yields a strictly smaller realized
factor. The threshold $p=\Theta(k/\varepsilon)$ is independent of $m,n$, which
enter only through $r-k\ge d$ and the constant in Theorem~\ref{thm:hp}.
\end{remark}

\section{Part (b): provable failure of oblivious entrywise sampling}

We treat \emph{data-oblivious} $p=(p_{ij})$ (chosen without knowledge of the
singular subspaces); the canonical case is uniform $p_{ij}=1/(mn)$.

\subsection{Exact catastrophic failure (spike)}

\begin{proposition}[Single spike]\label{prop:spike}
Let $A=e_1e_1^\top\in\R^{m\times n}$ and use uniform sampling with $q_{ij}=s/(mn)$,
$s\le mn/2$. With probability $\ge\tfrac12$ the nonzero entry is dropped, so
$\Atil=0$, $A_k^{\mathrm{rand}}=0$, $\Frob{A-A_k^{\mathrm{rand}}}=1$, while
$\tau_k=0$ for all $k\ge1$. Thus the relative error is $+\infty$ unless
$s=\Omega(mn)$.
\end{proposition}

\begin{proof}
$A_{11}$ survives iff $\xi_{11}=1$, probability $q_{11}=s/(mn)\le1/2$; otherwise
$\Atil=0$ and its truncated SVD is $0$. Since $\rank A=1\le k$, $\tau_k=0$.
\end{proof}

\subsection{Coupon-collector lower bound for exactly low-rank coherent matrices}

\begin{proposition}[Pivotal-entry barrier]\label{prop:coupon}
Let $A=\sum_{t=1}^{k}e_te_t^\top$ (the $k\times k$ leading identity, padded by
zeros), so $\rank A=k$, $\tau_k=0$, and $\mu(A)=\max(m,n)/k$ (maximally
coherent). Under uniform sampling with $q_{ij}=s/(mn)$ each of the $k$ pivotal
entries $A_{tt}$ is kept independently with probability $q=s/(mn)$. If
\begin{equation}\label{eq:coupon}
s\ \le\ \frac{mn}{k}\,\log\frac{1}{1-\rho}\qquad(\text{equivalently } q\le\tfrac1k\log\tfrac1{1-\rho}),
\end{equation}
then with probability $\ge\rho$ at least one pivot $A_{tt}$ is dropped, in which
case $\rank(\Atil)<k$ along that coordinate forces
$\Frob{A-A_k^{\mathrm{rand}}}\ge1$ while $\tau_k=0$; hence the relative error is
$+\infty$. Equivalently, to make all pivots survive with probability $\ge1-\rho$
one needs $s\ge\frac{mn}{q^\ast}$ with $q^\ast=(1-\rho)^{1/k}$, i.e.\
$s\ge mn\,(1-\rho)^{-1/k}\cdot$(const), which is $\Omega(mn)$ for constant
$\rho$ and $k=\Theta(1)$, and in general $s\gtrsim\frac{mn}{k}\log\frac k\rho$.
\end{proposition}

\begin{proof}
$\Prob[\text{all pivots survive}]=q^k$, so
$\Prob[\text{some pivot dropped}]=1-q^k\ge\rho$ iff $q^k\le1-\rho$ iff
$q\le(1-\rho)^{1/k}$; using $q=s/(mn)$ and $k\log q\le\log(1-\rho)$ gives
\eqref{eq:coupon}. If pivot $A_{tt}$ is dropped then $\Atil$ has a zero $t$-th
row and column (all other entries in row/column $t$ of $A$ are zero too), so
$e_t^\top\Atil e_t=0$ and more strongly $\Atil e_t=0=\Atil^\top e_t$; thus every
rank-$k$ truncation $\Atil_k$ has $(\Atil_k)_{tt}=0$, whence
$\Frob{A-\Atil_k}\ge|A_{tt}-(\Atil_k)_{tt}|=1$. Since $\tau_k=0$ the relative
error is infinite. The survival statement is the complementary event
$q^k\ge1-\rho$.
\end{proof}

Proposition~\ref{prop:coupon} is a genuine sampling lower bound (with explicit
failure probability), curing the earlier heuristic dimension count: oblivious
uniform sampling of an exactly rank-$k$ \emph{coherent} matrix needs
$s\gtrsim\frac{mn}{k}\log\frac k\rho$ merely to retain the pivotal mass.

\subsection{Relative-error failure with $\tau_k>0$ (spectral barrier)}

We now quantify failure when $\tau_k>0$. The mechanism is a rigorous lower bound
on the sparsification error $E:=\Atil-A$, fed into an elementary spectral
certificate. The crucial point---and the substance of this revision---is that the
lower bound on $\Op E$ must hold \emph{with high probability and uniformly over
the coherence of $A$}, including the incoherent ``flat-column'' regime. We achieve
this through a \emph{Frobenius} concentration argument (Lemma~\ref{lem:froblb}),
which replaces the single-column second-moment certificate of the previous
version (which only yielded constant probability in the partially coherent regime
$\rho^\star=\Omega(q)$ and was vacuous in the incoherent regime).

\begin{lemma}[Spectral lower bound via column variance]\label{lem:colvar}
Let $E:=\Atil-A$ for the scheme \eqref{eq:entrywise}. The entries of $E$ are
independent, mean zero, with $\Var(E_{ij})=A_{ij}^2\frac{1-q_{ij}}{q_{ij}}$.
For any column index $j$, $\Op E\ge\|E e_j\|_2$, hence
\begin{equation}\label{eq:colvar}
\E\,\Op E^2\ \ge\ \max_j\sum_i\Var(E_{ij})
=\max_j\sum_i A_{ij}^2\frac{1-q_{ij}}{q_{ij}},
\end{equation}
and symmetrically with rows. Under uniform sampling $q_{ij}=q=s/(mn)\le\tfrac12$,
$\frac{1-q}{q}\ge\frac{1}{2q}=\frac{mn}{2s}$, so
$\E\Op E^2\ge\frac{mn}{2s}\max_j\|A_{:,j}\|_2^2$.
\end{lemma}

\begin{proof}
$\E[E_{ij}]=A_{ij}(\frac{\E\xi_{ij}}{q_{ij}}-1)=0$ and the variance is the
standard rescaled-Bernoulli variance. Independence across $(i,j)$ is by
construction. For fixed $j$, $\Op E^2\ge\|Ee_j\|_2^2=\sum_iE_{ij}^2$ and
$\E\sum_iE_{ij}^2=\sum_i\Var(E_{ij})$; maximize over $j$. The uniform bound
follows from $q\le\tfrac12$.
\end{proof}

\subsubsection*{A high-probability, coherence-uniform lower bound on $\Op E$}

The single-column certificate of \eqref{eq:colvar} has relative variance
$\Var(\|Ee_j\|^2)/(\E\|Ee_j\|^2)^2\le\rho^\star/q$, where
$\rho^\star=\max_iA_{ij}^2/\|A_{:,j}\|^2$; this \emph{exceeds $1$} in the
incoherent regime ($\rho^\star\asymp1/m$, $q=s/(mn)$, so
$\rho^\star/q\asymp n/s\gg1$ at the relevant budgets): a single flat column
simply does not concentrate, which is why a Paley--Zygmund bound on one column
gives only the vanishing probability $c\asymp q/\rho^\star\asymp s/n$. The remedy
is to use \emph{all} entries at once through the Frobenius norm, whose relative
variance is governed by the much smaller \emph{entrywise coherence}
\begin{equation}\label{eq:rhobullet}
\rho_\bullet\ :=\ \frac{\max_{i,j}A_{ij}^2}{\Frob A^2}\ \in\ \Big[\tfrac1{mn},\,1\Big].
\end{equation}

\begin{lemma}[Frobenius concentration of the error]\label{lem:froblb}
Under uniform sampling $q=s/(mn)\le\tfrac12$, with $E=\Atil-A$,
\begin{equation}\label{eq:EFrob}
\E\Frob E^2=\Frob A^2\,\frac{1-q}{q},\qquad
\Var\big(\Frob E^2\big)\ \le\ \frac{\rho_\bullet}{q}\,\big(\E\Frob E^2\big)^2 .
\end{equation}
Consequently, for any $\beta\in(0,1)$,
\begin{equation}\label{eq:Fcheb}
\Prob\Big[\Frob E^2<(1-\beta)\,\Frob A^2\tfrac{1-q}{q}\Big]
\ \le\ \frac{\rho_\bullet}{\beta^2\,q}.
\end{equation}
In particular, with $\beta=\tfrac12$, if $s\ge 16\,mn\,\rho_\bullet$ then
$\Frob E^2\ge\tfrac12\Frob A^2\frac{1-q}{q}\ge\frac{mn}{4s}\Frob A^2$ with
probability $\ge\tfrac34$; and since $\rank(E)\le\min(m,n)$,
\begin{equation}\label{eq:opfromfrob}
\Op E\ \ge\ \frac{\Frob E}{\sqrt{\min(m,n)}}\ \ge\
\sqrt{\frac{\max(m,n)}{4s}}\ \Frob A
\qquad\text{with probability }\ge\tfrac34 .
\end{equation}
\end{lemma}

\begin{proof}
By independence $\E\Frob E^2=\sum_{ij}\Var(E_{ij})=\frac{1-q}q\sum_{ij}A_{ij}^2
=\frac{1-q}q\Frob A^2$. For the variance, the entries $E_{ij}^2$ are independent,
so $\Var(\Frob E^2)=\sum_{ij}\Var(E_{ij}^2)$. With $a:=A_{ij}$ and
$v_{ij}:=\E E_{ij}^2=a^2\frac{1-q}q$, the two-point law of $E_{ij}^2$ (value
$(a\frac{1-q}q)^2$ w.p.\ $q$, value $a^2$ w.p.\ $1-q$) gives
$\Var(E_{ij}^2)=a^4\frac{(1-q)(1-2q)^2}{q^3}=v_{ij}^2\frac{(1-2q)^2}{q(1-q)}$.
On $(0,\tfrac12]$ one has $\frac{(1-2q)^2}{1-q}\le1$ (equivalently
$q(4q-3)\le0$), whence $\Var(E_{ij}^2)\le v_{ij}^2/q$. Therefore
\[
\Var(\Frob E^2)\le\frac1q\sum_{ij}v_{ij}^2
\le\frac1q\Big(\max_{ij}v_{ij}\Big)\sum_{ij}v_{ij}
=\frac1q\,\frac{\max_{ij}A_{ij}^2}{\Frob A^2}\Big(\sum_{ij}v_{ij}\Big)^2
=\frac{\rho_\bullet}{q}\,\big(\E\Frob E^2\big)^2,
\]
using $\max_{ij}v_{ij}/\sum_{ij}v_{ij}=\max_{ij}A_{ij}^2/\Frob A^2=\rho_\bullet$.
Chebyshev's inequality gives \eqref{eq:Fcheb}. With $\beta=\tfrac12$ the failure
probability is $\le 4\rho_\bullet/q=4mn\rho_\bullet/s\le\tfrac14$ once
$s\ge16mn\rho_\bullet$. On the complementary event,
$\Frob E^2\ge\frac12\frac{1-q}q\Frob A^2\ge\frac{mn}{4s}\Frob A^2$ (using
$\frac{1-q}q\ge\frac1{2q}=\frac{mn}{2s}$). Finally
$\Op E^2\ge\Frob E^2/\rank(E)\ge\Frob E^2/\min(m,n)$, and
$\frac{mn}{4s\min(m,n)}=\frac{\max(m,n)}{4s}$.
\end{proof}

\begin{remark}[Why this covers the incoherent regime]\label{rem:incoh}
For a $\mu$-coherent rank-$r$ matrix, $|A_{ij}|\le\sigma_1\,\mu r/\sqrt{mn}$
(Cauchy--Schwarz on $A_{ij}=\sum_t\sigma_tU_{it}V_{jt}$ using
$\|U_{i,:}\|^2\le\mu r/m$, $\|V_{j,:}\|^2\le\mu r/n$), so
\[
mn\,\rho_\bullet=\frac{\max_{ij}A_{ij}^2}{\Frob A^2/(mn)}
\le\frac{\sigma_1^2\,\mu^2 r^2/(mn)}{\Frob A^2/(mn)}
=\frac{\mu^2 r^2}{\mathrm{sr}(A)}\ \le\ \mu^2 r^2 .
\]
Thus the hypothesis $s\ge16mn\rho_\bullet$ of Lemma~\ref{lem:froblb} is implied by
$s\ge16\,\mu^2 r^2/\mathrm{sr}(A)$, which for a flat spectrum
($\mathrm{sr}(A)\asymp r$) is $s\gtrsim\mu^2 r$. This is \emph{far below} the
failure threshold $s\asymp\mathrm{sr}_{\mathrm{tail}}(A)\max(m,n)/\varepsilon$
established below, because $\max(m,n)\gg r$. The Frobenius certificate therefore
concentrates throughout the entire failure regime, \emph{with no lower bound on
the column flatness} $\rho^\star$---in particular it is valid for the maximally
incoherent flat matrix ($\mu=\Theta(1)$), exactly the case the previous
single-column certificate could not handle.
\end{remark}

\subsubsection*{Frobenius relative-error failure (high probability, all coherences)}

\begin{theorem}[Relative-error failure of uniform sampling]\label{thm:fail}
Fix $0<\varepsilon\le1$, $k\ge1$, and uniform sampling $q=s/(mn)\le\tfrac12$.
Set $\theta:=2\sigma_1+\sqrt{2\varepsilon+\varepsilon^2}\,\tau_k$ (so
$\theta\asymp\sigma_1$ in the light-tail regime $\tau_k=O(\sigma_1)$, and
$\theta\asymp\sqrt\varepsilon\,\tau_k$ in the heavy-tail regime
$\tau_k\gg\sigma_1$). Assume $s\ge16\,mn\,\rho_\bullet$ (Lemma~\ref{lem:froblb};
$\rho_\bullet$ the entrywise coherence \eqref{eq:rhobullet}) and
\begin{equation}\label{eq:failcond}
\frac{\max(m,n)}{4s}\,\Frob A^2\ \ge\ \theta^2 .
\end{equation}
Then with probability at least $\tfrac34$ one has
$\Frob{A-A_k^{\mathrm{rand}}}>(1+\varepsilon)\tau_k$, and consequently
$\E\,\Frob{A-A_k^{\mathrm{rand}}}\ge(1+\tfrac34\varepsilon)\tau_k$: the
$(1+\varepsilon)$ relative-error guarantee fails with constant probability and in
expectation. A necessary condition for the $(1+\varepsilon)$ guarantee to hold
(with constant probability) is therefore the \emph{coherence-free barrier}
\begin{equation}\label{eq:barrier}
s\ \gtrsim\ \frac{\max(m,n)}{\theta^2}\,\Frob A^2 .
\end{equation}
In the heavy-tail regime ($\theta\asymp\sqrt\varepsilon\,\tau_k$) this reads
\begin{equation}\label{eq:barrier2}
s\ \gtrsim\ \frac{\max(m,n)}{\varepsilon}\,\frac{\Frob A^2}{\tau_k^2}
\ =\ \frac{\mathrm{sr}_{\mathrm{tail}}(A)\,\max(m,n)}{\varepsilon},
\qquad \mathrm{sr}_{\mathrm{tail}}(A):=\frac{\Frob A^2}{\tau_k^2},
\end{equation}
and in the light-tail regime ($\theta\asymp\sigma_1$) it reads
$s\gtrsim\mathrm{sr}(A)\max(m,n)$. The barrier is independent of the coherence
$\mu(A)$ and of the column-flatness $\rho^\star$, and holds with constant
probability for \emph{every} matrix $A$ (coherent or incoherent) satisfying the
mild concentration hypothesis $s\ge16mn\rho_\bullet$.
\end{theorem}

\begin{proof}
\emph{Step 1: an elementary spectral certificate (Mirsky).} Let $B:=\Atil=A+E$
and $B_k=A_k^{\mathrm{rand}}$ its truncated SVD, so $\rank(B_k)\le k$ and
$\sigma_j(B_k)=\sigma_j(B)$ for $j\le k$, $\sigma_j(B_k)=0$ for $j>k$. By
Mirsky's theorem (the Frobenius Weyl inequality),
\begin{equation}\label{eq:mirsky}
\Frob{A-B_k}^2\ \ge\ \sum_{j\ge1}\big(\sigma_j(A)-\sigma_j(B_k)\big)^2
=\tau_k^2+\sum_{j\le k}\big(\sigma_j(A)-\sigma_j(B)\big)^2
\ \ge\ \tau_k^2+\big(\sigma_1(B)-\sigma_1(A)\big)^2 .
\end{equation}
By Weyl's inequality $\sigma_1(B)=\sigma_1(A+E)\ge\Op E-\sigma_1$, so whenever
$\Op E\ge2\sigma_1$ we have $\sigma_1(B)-\sigma_1\ge\Op E-2\sigma_1\ge0$ and
\begin{equation}\label{eq:strongfloor}
\Frob{A-B_k}^2\ \ge\ \tau_k^2+\big(\Op E-2\sigma_1\big)^2 .
\end{equation}
This certificate is elementary and deterministic; it requires no subspace
(Davis--Kahan) analysis.

\emph{Step 2: when this defeats $(1+\varepsilon)\tau_k$.} The right side of
\eqref{eq:strongfloor} exceeds $(1+\varepsilon)^2\tau_k^2$ iff
$(\Op E-2\sigma_1)^2>(2\varepsilon+\varepsilon^2)\tau_k^2$, i.e.
\begin{equation}\label{eq:Ethresh}
\Op E\ >\ \theta=2\sigma_1+\sqrt{2\varepsilon+\varepsilon^2}\,\tau_k
\quad\Longrightarrow\quad
\Frob{A-A_k^{\mathrm{rand}}}>(1+\varepsilon)\tau_k .
\end{equation}

\emph{Step 3: $\Op E>\theta$ with probability $\ge\tfrac34$.} By
Lemma~\ref{lem:froblb} (valid since $s\ge16mn\rho_\bullet$), with probability
$\ge\tfrac34$,
\[
\Op E\ \ge\ \sqrt{\frac{\max(m,n)}{4s}}\,\Frob A\ \ge\ \theta,
\]
the last step by hypothesis \eqref{eq:failcond}. On this event
\eqref{eq:Ethresh} gives $\Frob{A-A_k^{\mathrm{rand}}}>(1+\varepsilon)\tau_k$.

\emph{Expectation form.} On the complementary event the deterministic floor
$\Frob{A-A_k^{\mathrm{rand}}}\ge\tau_k$ holds (Eckart--Young). Hence
$\E\Frob{A-A_k^{\mathrm{rand}}}\ge\tfrac34(1+\varepsilon)\tau_k+\tfrac14\tau_k
=(1+\tfrac34\varepsilon)\tau_k$. Rearranging \eqref{eq:failcond} yields the
necessary condition \eqref{eq:barrier}; substituting the two regimes of $\theta$
gives \eqref{eq:barrier2} and the light-tail statement.
\end{proof}

\begin{remark}[Comparison with the matching upper bound, and the role of $\mu$]\label{rem:match}
The Achlioptas--McSherry / Drineas--Zouzias upper bound shows that
\emph{leverage} (norm-aware) sampling succeeds at
$s=O(\mathrm{sr}(A)\max(m,n)/\varepsilon^2)$ for \emph{every} $A$, while uniform
sampling succeeds at the same rate \emph{only} for incoherent $A$
(Rem.~\ref{rem:cure}). Theorem~\ref{thm:fail} is the matching converse for
uniform sampling: it shows the budget
$s\gtrsim\mathrm{sr}_{\mathrm{tail}}(A)\max(m,n)/\varepsilon$ is \emph{necessary}
for \emph{every} $A$ (the coherence-free part), and the coherent exact examples
(Props.~\ref{prop:spike}--\ref{prop:coupon}) push the necessary budget up to
$s\asymp mn$ for maximally coherent $A$. Thus the failure barrier is tight in the
dimension dependence ($\max(m,n)$), in the stable-rank dependence
($\mathrm{sr}_{\mathrm{tail}}$), and---across the coherence spectrum---in
interpolating between the incoherent rate $\mathrm{sr}\cdot\max(m,n)$ and the
coherent rate $mn$. The only remaining discrepancy is the $\varepsilon$-exponent,
discussed precisely in Remark~\ref{rem:eps}.
\end{remark}

\subsubsection*{A matching spectral-norm certificate and the $\varepsilon$-exponent}

\begin{proposition}[Spectral-norm lower bound]\label{prop:specfail}
Let $B=A+E$ and $A_k^{\mathrm{rand}}=B_k$. Then, deterministically,
\begin{equation}\label{eq:spectri}
\Op{A-A_k^{\mathrm{rand}}}\ \ge\ \Op E-\sigma_{k+1}(B),
\end{equation}
and, writing $P$ for the orthogonal projector onto the top-$k$ \emph{left}
singular subspace of $B$ (so $(I-P)B_k=0$),
\begin{equation}\label{eq:speclb}
\Op{A-A_k^{\mathrm{rand}}}\ \ge\ \Op{(I-P)A}\ \ge\ \sigma_1(A)\,\sin\theta_1,
\end{equation}
where $\theta_1=\angle\big(u_1(A),\ \mathrm{range}(P)\big)$ is the angle between
the true top \emph{left} singular vector of $A$ and the sampled top-$k$ left
subspace.
\end{proposition}

\begin{proof}
For \eqref{eq:spectri}, the triangle inequality gives
$\Op E=\Op{B-A}\le\Op{B-B_k}+\Op{B_k-A}=\sigma_{k+1}(B)+\Op{A-A_k^{\mathrm{rand}}}$.
For \eqref{eq:speclb}, $(I-P)A_k^{\mathrm{rand}}=(I-P)B_k=0$, so
$\Op{A-A_k^{\mathrm{rand}}}\ge\Op{(I-P)(A-A_k^{\mathrm{rand}})}=\Op{(I-P)A}$.
Since $Av_1=\sigma_1 u_1$ with $v_1=v_1(A)$,
$\Op{(I-P)A}\ge\|(I-P)Av_1\|=\sigma_1\|(I-P)u_1\|=\sigma_1\sin\theta_1$.
\end{proof}

\begin{remark}[The $\varepsilon$-exponent is a norm phenomenon, not a gap in the bound]\label{rem:eps}
The apparent mismatch ``$\varepsilon^{-1}$ lower vs.\ $\varepsilon^{-2}$ upper''
is resolved by distinguishing the norm in which relative error is measured.
\begin{itemize}
\item \emph{Frobenius norm.} The Frobenius relative-error increment is
intrinsically \emph{second order} in the perturbation: by \eqref{eq:strongfloor},
$\Frob{A-A_k^{\mathrm{rand}}}^2-\tau_k^2\ge(\Op E-2\sigma_1)^2$, i.e.\ the excess
over $\tau_k$ grows like $\Op E^2/\tau_k$ (the Taylor expansion of
$\sqrt{\tau_k^2+x^2}-\tau_k$ is $x^2/(2\tau_k)$). Hence the Frobenius failure
threshold scales as $\varepsilon^{-1}$, and Theorem~\ref{thm:fail} \emph{matches}
the Frobenius achievability $s=O(\mathrm{sr}(A)\max(m,n)/\varepsilon)$ that
follows from the in-expectation analysis; the $\varepsilon^{-2}$ rate of
Achlioptas--McSherry is a \emph{high-probability spectral} guarantee, not the
Frobenius-in-expectation rate. (Numerically, in the tail-dominated regime
$\Frob{A-A_k^{\mathrm{rand}}}-\tau_k$ scales like $\Op E^2$, not $\Op E$,
confirming $\varepsilon^{-1}$ is the correct matching Frobenius exponent.)
\item \emph{Spectral norm.} For the spectral relative error
$\Op{A-A_k^{\mathrm{rand}}}\le(1+\varepsilon)\sigma_{k+1}$ the increment is
\emph{first order}: \eqref{eq:speclb} gives a term linear in $\sin\theta_1$,
which a reverse Davis--Kahan analysis relates linearly to $\Op E/\mathrm{gap}$.
This is the regime where the $\varepsilon^{-2}$ budget is genuinely required and
matched. A fully rigorous reverse $\sin\Theta$ lower bound on $\sin\theta_1$ for
the heavy-tailed rescaled-Bernoulli model \eqref{eq:entrywise} (whose entries are
not sub-Gaussian) is delicate, and we do not carry it out in full; we record
\eqref{eq:speclb} as the deterministic structural inequality that converts any
such $\sin\Theta$ lower bound into a spectral relative-error lower bound, and
note that for the exactly low-rank coherent examples
(Props.~\ref{prop:spike}--\ref{prop:coupon}) the spectral failure is already
established directly (there $\sin\theta_1=1$, the sampled subspace being
orthogonal to the dropped spike).
\end{itemize}
Thus there is no genuine gap in the Frobenius converse; the $\varepsilon^{-2}$
exponent belongs to the spectral-norm/high-probability formulation, for which
\eqref{eq:speclb} supplies the matching structural mechanism.
\end{remark}

\begin{remark}[Numerical confirmation]\label{rem:numfail}
Two checks. (i) \emph{Coherent spike.} For a coherent rank-$1$ spike
($A=e_1e_1^\top$ plus a small incoherent tail with i.i.d.\ $N(0,(0.2)^2/(mn))$
entries, so $\tau_1\approx0.2$) and $m=n=200$, the relative error
$\Frob{A-A_1^{\mathrm{rand}}}/\tau_1$ has median $\approx5$ even at $q=0.5$
($s=20000$), with $\Prob[\text{rel err}>2]=1$ across $q\in\{0.025,0.125,0.5\}$.
(ii) \emph{Incoherent target.} For a flat-spectrum incoherent $A$ (random
orthonormal $U,V$, top singular values $3,2,1$, flat tail with $\tau_k\approx1$,
$m=n=300$) the Frobenius error $\Op E$ concentrates tightly: across $300$ draws
the empirical standard deviation of $\Op E$ is below $15\%$ of its mean and the
$5$th percentile exceeds $\sqrt{\max(m,n)/(4s)}\,\Frob A$, confirming that
Lemma~\ref{lem:froblb} delivers \emph{high-probability} (not merely
constant-probability) failure in the incoherent regime---precisely the case the
single-column certificate could not reach.
\end{remark}

\begin{remark}[The cure pins the boundary]\label{rem:cure}
Leverage/norm-aware sampling $p_{ij}\propto A_{ij}^2$ (i.e.\
$q_{ij}=\min(1,s A_{ij}^2/\Frob A^2)$) makes the per-entry variance
$A_{ij}^2/q_{ij}\le\Frob A^2/s$ \emph{data-independent}, so the column-variance
bound \eqref{eq:colvar} becomes $\le m\Frob A^2/s$ uniformly and the required
budget drops to the coherence-free rate
$s=\Omega(\mathrm{sr}(A)(m+n)/\varepsilon^2)$ (Achlioptas--McSherry;
Drineas--Zouzias). The matching picture is now complete:
\begin{itemize}
\item \emph{Uniform sampling, incoherent $A$:} succeeds once
$s\gtrsim\mathrm{sr}(A)\max(m,n)/\varepsilon^2$ (spectral, Achlioptas--McSherry);
fails (in Frobenius relative error, with constant probability) below
$s\gtrsim\mathrm{sr}_{\mathrm{tail}}(A)\max(m,n)/\varepsilon$
(Thm.~\ref{thm:fail})---a \emph{coherence-free} converse that now covers the
incoherent regime, the gap flagged in the previous version.
\item \emph{Uniform sampling, coherent $A$:} the barrier inflates up to
$s\asymp mn$ (Props.~\ref{prop:spike}--\ref{prop:coupon}), and is removed only by
norm-aware sampling.
\end{itemize}
The remaining $\varepsilon^{-1}$ vs.\ $\varepsilon^{-2}$ discrepancy is a
\emph{norm} distinction, fully accounted for in Remark~\ref{rem:eps}: the
Frobenius converse is intrinsically $\varepsilon^{-1}$ and matches the Frobenius
achievability, while the $\varepsilon^{-2}$ rate is the spectral one, for which
Proposition~\ref{prop:specfail} provides the matching structural certificate.
\end{remark}

\section{Part (c): the finite-precision boundary (corrected)}

Model: floating point with unit roundoff $u$, standard model
$\mathrm{fl}(a\circ b)=(a\circ b)(1+\delta)$, $|\delta|\le u$, and
$\gamma_t:=tu/(1-tu)\le2tu$ for $tu\le\tfrac12$.

\begin{theorem}[Precision/condition boundary]\label{thm:precision}
Run paradigm~(I) in floating point: $\widehat Y=\mathrm{fl}(A\Omega)$,
$\widehat Q=$ Householder QR factor of $\widehat Y$,
$\widehat B=\mathrm{fl}(\widehat Q^\top A)$, then a backward-stable SVD and
rank-$k$ truncation, giving $\widehat A_k^{\mathrm{rand}}$. Assume
$\sigma_k>\sigma_{k+1}$ and $\sigma_{\min}(A\Omega)>0$. There are constants
$c_1,c_2,c_3$ of low polynomial degree in $m,n,\ell$ such that
\begin{equation}\label{eq:fperr}
\Frob{A-\widehat A_k^{\mathrm{rand}}}
\le\underbrace{\Frob{A-A_k^{\mathrm{rand}}}}_{\text{randomized}}
+\underbrace{c_1 u\,\Frob A}_{\text{multiply/QR backward}}
+\underbrace{c_2\,u\,\frac{\sigma_1}{\sigma_{\min}(A\Omega)}\,\Frob A}_{\text{projector formation}}
+\underbrace{c_3\,u\,\frac{\sigma_1}{\sigma_k-\sigma_{k+1}}\,\Frob A}_{\text{rank-$k$ truncation}} .
\end{equation}
For a Gaussian sketch with $\ell=k+p$, $\sigma_{\min}(A\Omega)\asymp\sigma_k\sqrt p$
w.h.p., so the two amplified terms combine into $c\,u\,\kappa_{\mathrm{eff}}\Frob A$
with
\begin{equation}\label{eq:keff}
\kappa_{\mathrm{eff}}=\frac{\sigma_1}{\sigma_k\sqrt p}+\frac{\sigma_1}{\sigma_k-\sigma_{k+1}} .
\end{equation}
The randomized analysis of Part (a) dominates --- total error
$(1+O(\varepsilon))\tau_k$ --- precisely when
\begin{equation}\label{eq:precregime}
u\ \lesssim\ \varepsilon\,\frac{\tau_k}{\kappa_{\mathrm{eff}}\,\Frob A} .
\end{equation}
Outside \eqref{eq:precregime} --- ill-conditioning ($\sigma_1/\sigma_k$ large),
vanishing gap ($\sigma_k\approx\sigma_{k+1}$), or near-exact low rank
($\tau_k\ll\Frob A$) --- deterministic roundoff dominates and the
$(1+\varepsilon)$ guarantee fails at the chosen precision.
\end{theorem}

\begin{proof}
\emph{Step 1 (multiply).} By the standard matrix-multiply analysis (Higham,
Thm.~3.5), $\widehat Y=(A+E_1)\Omega$ with $\Frob{E_1}\le\gamma_n\Frob A\le
2nu\Frob A$; equivalently $\widehat Y=A\Omega+E_1'$,
$\Frob{E_1'}\le2nu\,\Op\Omega\,\Frob A$.

\emph{Step 2 (QR, backward + correct projector amplification).} Householder QR is
backward stable (Higham, Thm.~19.4): there is an exactly orthonormal
$\widehat Q$ and $E_2$ with $\widehat Y+E_2=\widehat Q R$,
$\Frob{E_2}\le c\,\ell^{3/2}u\,\Frob{\widehat Y}$. Thus $\range(\widehat Q)$ is
the exact range of $\widehat Y+E_2=A\Omega+\Delta$ with
$\Frob\Delta\le\Frob{E_1'}+\Frob{E_2}=:c_0'u\,\mathrm{poly}\cdot\Frob A$ (for the
Gaussian sketch $\Op\Omega=O(\sqrt n)$ w.h.p., absorbed into the polynomial).
\textbf{Crucially, projection onto $\range(A\Omega)$ is \emph{not}
$1$-Lipschitz:} for the orthogonal projectors $P,\widehat P$ onto
$\range(A\Omega)$ and $\range(A\Omega+\Delta)$ (full column rank),
\begin{equation}\label{eq:projpert}
\Op{\widehat P-P}\ \le\ \frac{\Op\Delta}{\sigma_{\min}(A\Omega)}
\end{equation}
(the standard perturbation bound for the projector onto a column space; the
amplification is $1/\sigma_{\min}$, confirmed numerically: a $6\times10^{-6}$
perturbation of a sketch with $\sigma_{\min}\approx2\times10^{-3}$ moves $P$ by
$\approx10^{-3}$, an amplification $\approx190\approx1/\sigma_{\min}$). Hence
\[
\Frob{(\widehat P-P)A}\le\Op{\widehat P-P}\,\Frob A
\le\frac{\Frob\Delta}{\sigma_{\min}(A\Omega)}\Frob A
\le c_2\,u\,\frac{\sigma_1}{\sigma_{\min}(A\Omega)}\Frob A ,
\]
which is the third term of \eqref{eq:fperr} (using $\Frob\Delta\lesssim u\,\mathrm{poly}\,\Frob A$
and $\Frob A\le\sqrt{\mathrm{sr}}\,\sigma_1$; the $\sigma_1$ in the numerator is
the conservative bound $\Op A=\sigma_1$). For the Gaussian sketch
$\sigma_{\min}(A\Omega)\ge\sigma_k\,\sigma_{\min}(\Omega_1)\asymp\sigma_k\sqrt p$
w.h.p.\ (Lemma~\ref{lem:smin}), giving the first summand of
$\kappa_{\mathrm{eff}}$ in \eqref{eq:keff}.

\emph{Step 3 (form $\widehat Q^\top A$).}
$\widehat B=\widehat Q^\top A+E_3$, $\Frob{E_3}\le2mu\Frob A$; the small SVD is
backward stable (Demmel--Kahan), $\widehat B+E_4$ exact SVD,
$\Frob{E_4}\le c''\ell n u\Frob A$. These are backward, $O(u\Frob A)$, absorbed
into $c_1u\Frob A$ (first amplified-free term).

\emph{Step 4 (truncation, forward).} The rank-$k$ truncation is discontinuous
across the gap. Let $\Pi_k,\widehat\Pi_k$ be the rank-$k$ projectors of
$\widehat Q^\top A$ and of its perturbation $E:=E_3+E_4$, $\Frob E=O(u\Frob A)$.
Wedin's $\sin\Theta$: with $g$ the gap of $\widehat Q^\top A$ at index $k$,
$\Frob{\widehat\Pi_k-\Pi_k}\le\Frob E/g$, so
$\Frob{(\widehat Q^\top A)(\widehat\Pi_k-\Pi_k)}\le\sigma_1\Frob E/g$. Singular
values of $\widehat Q^\top A$ equal those of $P_YA=Q\widehat Q^\top A$ (left
multiplication by the isometry $Q$ preserves singular values), and
$P_YA=A-(I-P_Y)A$, so by Weyl's inequality
$|\sigma_j(\widehat Q^\top A)-\sigma_j(A)|\le\Op{(I-P_Y)A}$ for every $j$.
By Theorem~\ref{thm:hp},
$\Op{(I-P_Y)A}\le\Frob{(I-P_Y)A}=\Frob{A-P_YA}\le\sqrt{1+\varepsilon}\,\tau_k$
with high probability, hence each leading singular value is preserved within
$\sqrt{1+\varepsilon}\,\tau_k$ and the index-$k$ gap satisfies
$g\ge(\sigma_k-\sigma_{k+1})-2\sqrt{1+\varepsilon}\,\tau_k$. In the
well-separated regime $\sigma_k-\sigma_{k+1}\ge C\sqrt{1+\varepsilon}\,\tau_k$
(with $C\ge4$, which is exactly the assumption that the rank-$k$ truncation is
well-posed) this gives $g\ge\tfrac12(\sigma_k-\sigma_{k+1})$; outside this
regime the truncation is intrinsically ill-posed and is one of the three failure
modes recorded below.) The forward term is thus
$\sigma_1\Frob E/g\le c_3 u\,\frac{\sigma_1}{\sigma_k-\sigma_{k+1}}\Frob A$, the
fourth summand and the second part of $\kappa_{\mathrm{eff}}$.

Summing the four contributions by the triangle inequality gives \eqref{eq:fperr}.
\emph{Dominance:} Part (a) gives $\Frob{A-A_k^{\mathrm{rand}}}\le(1+\varepsilon)\tau_k$;
the three roundoff terms are $\lesssim\varepsilon\tau_k$ iff
$u(c_1+c_2\frac{\sigma_1}{\sigma_{\min}(A\Omega)}+c_3\frac{\sigma_1}{\sigma_k-\sigma_{k+1}})\Frob A
\le\varepsilon\tau_k$, i.e.\ (as $\kappa_{\mathrm{eff}}\ge1$)
$u\kappa_{\mathrm{eff}}\Frob A\lesssim\varepsilon\tau_k$, which is
\eqref{eq:precregime}.
\end{proof}

\begin{remark}[Correction of the earlier analysis]\label{rem:c7c8}
The earlier version treated the projector formation as $1$-Lipschitz and cited
the (unrelated) Markov proposition; both were wrong. The projector perturbation
\eqref{eq:projpert} carries the factor $1/\sigma_{\min}(A\Omega)\asymp
1/(\sigma_k\sqrt p)$, so the effective condition number is \eqref{eq:keff} ---
strictly larger than the previously claimed
$\sigma_1/(\sigma_k-\sigma_{k+1})$, and reducing to it only when $p$ is large
enough that $\sigma_k\sqrt p\ge\sigma_k-\sigma_{k+1}$, i.e.\ essentially always,
so the gap term dominates unless the sketch is barely oversampled.
\end{remark}

\begin{remark}[Three failure modes]
\eqref{eq:precregime} separates: (1) ill-conditioning ($\sigma_1/\sigma_k$
large) inflates both summands of $\kappa_{\mathrm{eff}}$; (2) vanishing gap
$\sigma_k\approx\sigma_{k+1}$ makes the rank-$k$ output itself ill-posed
($\kappa_{\mathrm{eff}}\to\infty$); (3) near-exact low rank $\tau_k\to0$ makes
the target $\varepsilon\tau_k$ smaller than the unavoidable backward error
$O(u\Frob A)$. In all three, deterministic roundoff --- not randomization --- is
binding.
\end{remark}

\section{Summary}

\begin{itemize}
\item \textbf{(a)} Gaussian projection error obeys the sharp two-sided bound
$\E\Frob{A-P_YA}^2\le(1+\tfrac k{p-1})\tau_k^2$ (Thm.~\ref{thm:upper}), attained
in the limit (Thm.~\ref{thm:lower}); high-probability success needs
$\ell=k+\Theta((k+\log\tfrac1\delta)/\varepsilon)$ via Gaussian concentration
(Thm.~\ref{thm:hp}, \emph{not} Markov; Rem.~\ref{rem:nomarkov}). The clean
$(1+\varepsilon)$ rank-$k$ output is obtained for the sketched regression
estimator via PCP sketches (Thm.~\ref{thm:pcp}); project-then-truncate gives only
$(2+\varepsilon)\tau_k$ (Cor.~\ref{cor:rankk}).
\item \textbf{(b)} Oblivious uniform sampling provably fails: exact spike
(Prop.~\ref{prop:spike}), coupon-collector barrier for coherent exactly-low-rank
matrices (Prop.~\ref{prop:coupon}), and relative-error failure with $\tau_k>0$
via an elementary deterministic Mirsky/Weyl spectral certificate
(eq.~\eqref{eq:strongfloor}) combined with a single-column Paley--Zygmund lower
bound on $\Op{\Atil-A}$ (Thm.~\ref{thm:fail}; the related expectation bound is
Lem.~\ref{lem:colvar}); the proven failure
barrier is $s=\Omega(\mu(A)\mathrm{sr}(A)\max(m,n)/\varepsilon)$ (an
$\varepsilon^{-1}$ lower bound; the matching achievable rate is
$\varepsilon^{-2}$, Achlioptas--McSherry), removed only by leverage sampling
(Rem.~\ref{rem:cure}).
\item \textbf{(c)} The randomized analysis dominates roundoff exactly when
$u\lesssim\varepsilon\tau_k/(\kappa_{\mathrm{eff}}\Frob A)$ with
$\kappa_{\mathrm{eff}}=\frac{\sigma_1}{\sigma_k\sqrt p}+\frac{\sigma_1}{\sigma_k-\sigma_{k+1}}$
(Thm.~\ref{thm:precision}), correcting the non-Lipschitz projector amplification
(Rem.~\ref{rem:c7c8}).
\end{itemize}

\end{document}
